%Version 3.1 December 2024
% See section 11 of the User Manual for version history
%
%%%%%%%%%%%%%%%%%%%%%%%%%%%%%%%%%%%%%%%%%%%%%%%%%%%%%%%%%%%%%%%%%%%%%%
%%                                                                 %%
%% Please do not use \input{...} to include other tex files.       %%
%% Submit your LaTeX manuscript as one .tex document.              %%
%%                                                                 %%
%% All additional figures and files should be attached             %%
%% separately and not embedded in the \TeX\ document itself.       %%
%%                                                                 %%
%%%%%%%%%%%%%%%%%%%%%%%%%%%%%%%%%%%%%%%%%%%%%%%%%%%%%%%%%%%%%%%%%%%%%

%%\documentclass[referee,sn-basic]{sn-jnl}% referee option is meant for double line spacing

%%=======================================================%%
%% to print line numbers in the margin use lineno option %%
%%=======================================================%%

%%\documentclass[lineno,pdflatex,sn-basic]{sn-jnl}% Basic Springer Nature Reference Style/Chemistry Reference Style

%%=========================================================================================%%
%% the documentclass is set to pdflatex as default. You can delete it if not appropriate.  %%
%%=========================================================================================%%

%%\documentclass[sn-basic]{sn-jnl}% Basic Springer Nature Reference Style/Chemistry Reference Style

%%Note: the following reference styles support Namedate and Numbered referencing. By default the style follows the most common style. To switch between the options you can add or remove “Numbered” in the optional parenthesis. 
%%The option is available for: sn-basic.bst, sn-chicago.bst%  
 
%%\documentclass[pdflatex,sn-nature]{sn-jnl}% Style for submissions to Nature Portfolio journals
%%\documentclass[pdflatex,sn-basic]{sn-jnl}% Basic Springer Nature Reference Style/Chemistry Reference Style
%\documentclass[pdflatex,sn-mathphys-num]{sn-jnl}
\documentclass[sn-mathphys-num]{sn-jnl}

% Math and Physical Sciences Numbered Reference Style
%%\documentclass[pdflatex,sn-mathphys-ay]{sn-jnl}% Math and Physical Sciences Author Year Reference Style
%%\documentclass[pdflatex,sn-aps]{sn-jnl}% American Physical Society (APS) Reference Style
%%\documentclass[pdflatex,sn-vancouver-num]{sn-jnl}% Vancouver Numbered Reference Style
%%\documentclass[pdflatex,sn-vancouver-ay]{sn-jnl}% Vancouver Author Year Reference Style
%%\documentclass[pdflatex,sn-apa]{sn-jnl}% APA Reference Style
%%\documentclass[pdflatex,sn-chicago]{sn-jnl}% Chicago-based Humanities Reference Style

%%%% Standard Packages
%%<additional latex packages if required can be included here>

\usepackage{graphicx}%
\usepackage{multirow}%
\usepackage{amsmath,amssymb,amsfonts}%
\DeclareMathOperator{\sign}{sign}
\usepackage{amsthm}%
\usepackage{mathrsfs}%
\usepackage[title]{appendix}%
\usepackage{xcolor}%
\usepackage{textcomp}%
\usepackage{manyfoot}%
\usepackage{booktabs}%
\usepackage{algorithm}%
\usepackage{algorithmicx}%
\usepackage{algpseudocode}%
\usepackage{listings}%
%%%%
\usepackage{txfonts}
\usepackage[utf8]{inputenc}
\usepackage{makecell}
\usepackage{xcolor}  \usepackage{soul}
\newtheorem{assumption}{Assumption}

%%%%%=============================================================================%%%%
%%%%  Remarks: This template is provided to aid authors with the preparation
%%%%  of original research articles intended for submission to journals published 
%%%%  by Springer Nature. The guidance has been prepared in partnership with 
%%%%  production teams to conform to Springer Nature technical requirements. 
%%%%  Editorial and presentation requirements differ among journal portfolios and 
%%%%  research disciplines. You may find sections in this template are irrelevant 
%%%%  to your work and are empowered to omit any such section if allowed by the 
%%%%  journal you intend to submit to. The submission guidelines and policies 
%%%%  of the journal take precedence. A detailed User Manual is available in the 
%%%%  template package for technical guidance.
%%%%%=============================================================================%%%%

%% as per the requirement new theorem styles can be included as shown below
\theoremstyle{thmstyleone}%
\newtheorem{theorem}{Theorem}%  meant for continuous numbers
%%\newtheorem{theorem}{Theorem}[section]% meant for sectionwise numbers
%% optional argument [theorem] produces theorem numbering sequence instead of independent numbers for Proposition
\newtheorem{proposition}[theorem]{Proposition}% 
%%\newtheorem{proposition}{Proposition}% to get separate numbers for theorem and proposition etc.

\theoremstyle{thmstyletwo}%
\newtheorem{example}{Example}%
\newtheorem{remark}{Remark}%

\theoremstyle{thmstylethree}%
\newtheorem{definition}{Definition}%

\raggedbottom
%%\unnumbered% uncomment this for unnumbered level heads

\begin{document}

\title[Article Title]{Robust Adaptive Control of Underactuated UAV-Suspended Payload Systems: A Symmetry Approach for Last-Mile Delivery}

%%=============================================================%%
%% GivenName	-> \fnm{Joergen W.}
%% Particle	-> \spfx{van der} -> surname prefix
%% FamilyName	-> \sur{Ploeg}
%% Suffix	-> \sfx{IV}
%% \author*[1,2]{\fnm{Joergen W.} \spfx{van der} \sur{Ploeg} 
%%  \sfx{IV}}\email{iauthor@gmail.com}
%%=============================================================%%
\author*[1]{\fnm{Souad} \sur{Kamel}}\email{skamel@uj.edu.sa}\equalcont{}

\author[2]{\fnm{Claudio} \sur{Rosales}}\email{crosales@inaut.unsj.edu.ar}\equalcont{}

\author[3]{\fnm{Sahbi} \sur{Boubaker}}\email{sboubaker@uj.edu.sa}\equalcont{}


\author[3]{\fnm{Faisal S.} \sur{Alsubaei}}\email{fsalsubaei@uj.edu.sa}\equalcont{}


\author[4]{\fnm{Eduardo Oliverira} \sur{Freire}}\email{efreire@academico.ufs.br}\equalcont{}


\author[2]{\fnm{Francisco G.} \sur{Rossomando}}\email{frosoma@inaut.unsj.edu.ar}\equalcont{}

\affil*[1]{\orgdiv{Department of Computer and Network Engineering}, \orgname{College of Computer Science and Engineering, University of Jeddah}, \orgaddress{\city{Jeddah}, \postcode{23890}, \state{Jeddah}, \country{Saudi Arabia}}}

\affil[2]{\orgdiv{Instituto de Automatica}, \orgname{UNSJ-CONICET}, \orgaddress{ \city{San Juan}, \postcode{CP5400}, \state{San Juan}, \country{Argentina}}}

\affil[3]{\orgdiv{Department of Cybersecurity}, \orgname{College of Computer Science and Engineering, University of Jeddah}, \orgaddress{\city{Jeddah}, \postcode{21959}, \state{Jeddah}, \country{Saudi Arabia}}}

\affil[4]{\orgdiv{Departamento de Engenharia Elétrica}, \orgname{Universidade Federal de Sergipe}, \orgaddress{\city{São Cristóvão/SE}, \postcode{49100-000}, \state{Sergipe}, \country{Brasil}}}


%%==================================%%
%% Sample for unstructured abstract %%
%%==================================%%

\abstract{This work addresses the computational modeling, simulation, and design of a composite adaptive control algorithm for last-mile delivery (LMD) using quadcopters (UAVs), which carry a payload. This work explores the interplay between the inherent asymmetry of coupled UAV-payload dynamics and the restoration of performance symmetry through a robust control approach. The system under study consists of the quadrotor, the cable, and the payload, where these elements are modeled as a subactuated dynamic system in which the payload must follow a desired trajectory and the quadrotor acts as the actuation platform. Taking advantage of the geometric and structural symmetries inherent in the pendulum-type configuration, a cascade control strategy is proposed that utilizes feedback linearization at both the kinematic and dynamic levels. To dampen oscillations and compensate for parametric uncertainties, a robust adaptive compensation algorithm based on the super-twisting algorithm (STA) is integrated. The proposal uses a composite tuning mechanism that effectively bridges the gap between the model’s nominal parameters and real-world dynamics. This strategy ensures computational efficiency and robustness during control. Stability analysis was performed using Lyapunov theory, rigorously demonstrating the asymptotic convergence of control errors and ensuring the stability of the closed-loop system. The numerical simulations presented for a quadrotor carrying a suspended payload validate the proposed control architecture, demonstrating that the proposal is effective in rejecting disturbances and ensuring precise trajectory tracking, making it ideal for high-precision UAV applications.}


%%================================%%
%% Sample for structured abstract %%
%%================================%%

% \abstract{\textbf{Purpose:} The abstract serves both as a general introduction to the topic and as a brief, non-technical summary of the main results and their implications. The abstract must not include subheadings (unless expressly permitted in the journal's Instructions to Authors), equations or citations. As a guide the abstract should not exceed 200 words. Most journals do not set a hard limit however authors are advised to check the author instructions for the journal they are submitting to.
% 
% \textbf{Methods:} The abstract serves both as a general introduction to the topic and as a brief, non-technical summary of the main results and their implications. The abstract must not include subheadings (unless expressly permitted in the journal's Instructions to Authors), equations or citations. As a guide the abstract should not exceed 200 words. Most journals do not set a hard limit however authors are advised to check the author instructions for the journal they are submitting to.
% 
% \textbf{Results:} The abstract serves both as a general introduction to the topic and as a brief, non-technical summary of the main results and their implications. The abstract must not include subheadings (unless expressly permitted in the journal's Instructions to Authors), equations or citations. As a guide the abstract should not exceed 200 words. Most journals do not set a hard limit however authors are advised to check the author instructions for the journal they are submitting to.
% 
% \textbf{Conclusion:} The abstract serves both as a general introduction to the topic and as a brief, non-technical summary of the main results and their implications. The abstract must not include subheadings (unless expressly permitted in the journal's Instructions to Authors), equations or citations. As a guide the abstract should not exceed 200 words. Most journals do not set a hard limit however authors are advised to check the author instructions for the journal they are submitting to.}

\keywords{Quadrotor UAV, last-mile delivery (LMD), transportation, adaptive control, trajectory tracking, STA, robust control}

%%\pacs[JEL Classification]{D8, H51}

%%\pacs[MSC Classification]{35A01, 65L10, 65L12, 65L20, 65L70}


\maketitle

\section{Introduction}

\subsection{State of the art and Control of Cable-Suspended Payload Transport}
Cable-suspended payload transport via unmanned aerial vehicles (UAVs) is a cornerstone of modern aerial logistics, infrastructure inspection, and emergency response in complex environments. Despite its operational potential, high-precision transport using a single UAV remains an open challenge. The system is fundamentally underactuated: the inclusion of a suspended load introduces oscillatory degrees of freedom that cannot be directly controlled, as the number of independent control inputs does not increase with the payload’s dimensionality. In this coupled quadcopter–cable–payload system, payload oscillations and external factors, such as wind gusts, act as primary sources of dynamic asymmetry that break the system’s natural equilibrium. These asymmetric perturbations not only degrade tracking accuracy but also introduce directional biases in the system's response, exacerbated by the nonlinear coupling \cite{Estevez2024,Wang2023,Yu2024,Mohammadi2021}.
Two configurations are commonly considered. Cooperative transport with several quadcopters distributes the load over more actuators and can improve transport performance \cite{Costantini2024,sun2025agile,di2026modeling,han2025sliding}, where these systems  require advanced, nonlinear control techniques to reject
disturbances \cite{han2022controller} and avoid obstacles \cite{aliyu2022control}, whereas single-UAV transport keeps the hardware, sensing, and coordination requirements minimal. Since the cooperative case and its drawbacks are discussed later in this section, the present work focuses on the single-vehicle configuration, in which one higher-performance vehicle must carry the suspended load along the desired trajectory. In much of the literature, the suspended load is treated as a disturbance or an undesirable effect to be rejected.
For single-vehicle transport, where quadcopters today play a significant role in logistics and rescue operations, most control designs are built upon a dynamic model of the coupled quadcopter--payload--cable system. Goodarzi et al. (2014) \citep{goodarzi2014geometric} proposed a geometric controller on the Lie group SE(3) that avoids the singularities of Euler-angle parameterizations. Its performance, however, is tied to an accurate full-dynamics model and full-state feedback, so that wind gusts and parametric uncertainty, explicitly acknowledged by the authors,  directly degrade tracking, while the scheme provides no online mechanism to estimate or compensate for unmodeled effects. To improve disturbance handling, Qian and Liu (2019) \citep{qian2019path} and Barikbin and Fakharian (2019) \citep{barikbin2019trajectory} developed robust controllers that preserve tracking accuracy under bounded disturbances. Nevertheless, robustness is typically obtained through conservative worst-case gains tuned around fixed operating points, which limits performance under time-varying disturbances and, as the authors themselves note, leaves part of the load swing uncompensated. Passivity-based strategies, such as the energy-shaping controller of Guerrero et al. (2015) \citep{guerrero2015passivity},  exploit the system's energy structure to attenuate load oscillations. This approach, however, requires an accurate energy and parameter model and generally yields only passive asymptotic damping, whose convergence is slow against persistent excitation. Sliding-mode control, used by Kusznir and Smoczek (2020) \citep{kusznir2020sliding}, is well suited to rejecting matched uncertainties and reducing payload vibration. Yet conventional first-order implementations act at the UAV position or attitude level, induce chattering that stresses the actuators, and rely on prior knowledge of the disturbance bounds. Finally, optimization- and planning-based methods, optimal trajectory tracking in Akhtar et al. (2022) \citep{akhtar2022path} and swing-free trajectory generation in Palunko et al. (2012) \cite{palunko2012trajectory}, reduce oscillations by shaping the reference itself. Their largely feedforward nature, however, makes them sensitive to model mismatch and unanticipated disturbances, and the associated optimization introduces a computational burden that is difficult to sustain on lightweight embedded platforms.
%%%%%
Beyond the single-vehicle case, transporting loads that exceed the capacity of one drone has motivated cooperative systems (Maza et al., 2010 \cite{maza2009multi}; Michael et al., 2011 \cite{michael2011cooperative}), in which coordination and joint manipulation are essential. These gains, however, come at the cost of additional hardware, distributed sensing, and inter-agent coordination, which restrict their applicability in the low-cost, single-UAV missions targeted in this work. A related line refines the physical realism of the model: cable-tension management to prevent the instantaneous loss of controllability under slack cables (Delbene, 2024) \cite{delbene2024cables}, geometric control accounting for cable elasticity (Goodman et al., 2023) \cite{goodman2023geometric}, and quantification of the tether effect on UAV stability through Lyapunov exponents (Tang et al., 2024) \cite{tang2024analyzing}. Valuable as they are, these contributions increase the model fidelity and state-estimation requirements, and they remain primarily analysis- or modeling-oriented rather than delivering a lightweight feedback law for real-time transport. A second group manages uncertainty through intelligent or adaptive compensation, including optimized neural PID controllers under unknown dynamics and external disturbances (Siwek et al., 2024) \cite{siwek2024application} and robust architectures tolerant to unknown payload mass (Outeiro et al., 2023) \cite{outeiro2023control}. Learning-based schemes, however, depend on representative training data and seldom provide formal stability guarantees, whereas mass-robust designs typically address a single uncertain parameter rather than the combined effect of parametric drift, unmodeled dynamics, and external disturbances. A third group operates at the planning and coordination level: two-time-scale swing minimization (Costantini et al., 2025) \cite{costantini2025two}, distributed trajectory optimization for scalability (Jackson et al., 2020) \cite{jackson2020scalable}, safe formation transport via ``virtual tubes'' (Liu et al., 2026) \cite{liu2026virtual}, and SQP-based tension allocation for energy efficiency and slack avoidance (Vazquez et al., 2025) \citep{vazquez:hal-05485942}. These methods presuppose either a clear time-scale separation, a multi-UAV formation, or the online solution of constrained optimization problems, all of which raise the computational and structural demands beyond what a single embedded platform can typically afford. Finally, perception- and orientation-aware controllers that decouple translational motion from yaw to preserve situational awareness during transport (Masuda and Sekiyama, 2026) \cite{masuda2026geometric} target viewpoint regulation rather than payload-swing reduction, and thus pursue an objective that is complementary to, but distinct from, the one addressed here.

Recent work on cable-suspended transport can be broadly grouped into three main directions. The first includes fully coupled nonlinear dynamic controllers, such as geometric, passivity-based, or model-based schemes derived from complete Newton--Euler or Lagrangian descriptions of the UAV--payload system \cite{Estevez2024,Costantini2024}. These approaches can achieve high accuracy when the model is reliable, but they usually require detailed knowledge of payload mass, cable geometry, and tension conditions, as well as relatively rich state information and nontrivial computational effort. The second direction comprises robust and adaptive sliding-mode-based strategies, including higher-order and super-twisting algorithms \cite{Wang2023,Mohammadi2021,rao2024adaptive}. These methods improve disturbance rejection and robustness, but in much of the literature, the sliding-mode action is applied at the position or attitude level, with limited integration into a structured payload-aware reference generation layer and only partial experimental validation for real suspended-load transport. In a recent study by Budak et al. (2026) \cite{budak2026hybrid}, a Linear Quadratic Regulator (LQR) with Sliding Mode Control (SMC) and Super-Twisting (STSMC) is used to improve transient performance in a 3-degree-of-freedom levitation system without a suspended load. The methodology uses the Big Bang–Big Crunch (BB–BC) optimization algorithm to automatically tune the controller parameters, thereby achieving an optimal balance between tracking accuracy and chattering suppression.


The third direction includes optimization-based and predictive methods, such as swing-aware trajectory planning and nonlinear MPC \cite{Yu2024,Estevez2024}. Their main advantage is the explicit handling of constraints and swing reduction at the planning stage, but they generally rely on accurate models and higher computational resources, which may hinder deployment on lightweight embedded platforms.
Recent advancements in agile aerial manipulation, such as the cinodynamic planning proposed by Sun et al. (2025) \cite{sun2025agile}, have pushed the boundaries of maneuverability. However, there is still a gap in providing lightweight, adaptive solutions that handle the asymmetric perturbations of wind and uncertain payload masses without the high computational cost of Nonlinear MPC. This work addresses this by utilizing an Adaptive Super-Twisting Algorithm (A-STA) that ensures tracking invariance under these conditions.

\subsection{Research gap and study contributions}

Although the reviewed studies have achieved notable advances in the control of UAV cable-suspended payload transport systems, they also highlight several remaining practical limitations. In particular, existing approaches rarely offer a lightweight control architecture capable of simultaneously:
(i) defining the control objective in payload space, (ii) separating payload-aware kinematic reference generation from robust inner-loop tracking, (iii) preserving robustness under uncertainty and disturbances, (iv) remaining suitable for implementation with limited sensing and computation while being supported by experimental validation. This gap is especially relevant in single-UAV transport, where the controller must reduce swing without relying on cooperative actuation or computationally expensive optimization.


To address this problem, this paper proposes a composite payload-centric control architecture composed of two complementary layers. The outer layer is a kinematic reference generator that shapes the UAV velocity command from payload tracking and cable-motion information in order to reduce oscillations during transport. The inner layer is an adaptive robust velocity-tracking controller based on a super-twisting sliding-mode structure, designed to realize the commanded motion in the presence of uncertainty and disturbances. In contrast with purely UAV-centric schemes, the proposed design explicitly prioritizes payload motion. In contrast with fully coupled model-heavy approaches, it adopts a modular structure with lower implementation complexity. In contrast with optimization-based methods, it does not require solving online constrained optimization problems. Finally, compared with conventional sliding-mode implementations reported in the literature, the proposed scheme integrates super-twisting robustness within a payload-aware cascaded architecture and evaluates the resulting controller through high-fidelity simulation experiments.

% REV1-C1/C3: Reforzar esta transición. Conviene declarar en forma inequívoca el foco del artículo: no es optimización logística end-to-end, sino control de transporte con carga suspendida dentro de un framework LMD. Agregar una frase tipo "The upper-level modules are not part of the controller synthesis; they provide feasible payload references and operational constraints."
% REV2-C1/C4: Revisar la frase "validates ... through flight experiments"; actualmente los resultados presentados son de simulación. Además, la afirmación comparativa con métodos convencionales debe respaldarse con una tabla/experimento contra controladores existentes.

The main contributions of this paper are as follows:

\begin{itemize}

 \item \textbf{Payload-Oriented Kinematic Control Design:} We introduce an external-loop kinematic formulation that shifts the control objective from the UAV chassis to the payload. By exploiting the kinematic relationship between the cable’s attitude and the quadcopter’s velocity, we derive oscillation compensation terms that act on the symmetry of the coupled system, mitigating residual roll by managing the velocity references.
 \item \textbf{Restore Tracking Symmetry:} A key original contribution of this architecture is its ability to restore tracking symmetry. By explicitly accounting for the payload's state in the kinematic layer, the controller ensures that tracking errors remain uniform and balanced across all translational axes, regardless of the asymmetric nature of external disturbances like oblique wind or pendulum swing.

 \item \textbf{Robust and Adaptive Cascaded Tracking Architecture:} We develop a hierarchical control structure that decouples the payload’s kinematics from the quadrotor’s internal dynamics. The inner loop employs an adaptive super-twisting speed controller, which ensures robustness against parametric uncertainty and external disturbances without degrading the transient response, thereby preserving system stability under varying operating conditions.

\item\textbf{Structured Control Paradigm vs. Conventional Approaches:} Our methodology marks a turning point in the existing literature:

\begin{itemize}
\item[$\ast$] Compared to coupled nonlinear controllers, our solution reduces the reliance on an exact dynamic model, thereby increasing resilience to modeling errors.

\item[$\ast$] Unlike heavy optimization solvers used in NMPC, our modular architecture and analytical adaptation laws allow implementation on low-power embedded hardware, which is a critical requirement for maintaining control symmetry and reliability in real-world last-mile delivery drones.

\item[$\ast$] Unlike standard sliding mode control (SMC) implementations, our implementation prioritizes payload stability, preventing corrective actions from the inner loop from generating unwanted oscillations in the suspension.
\end{itemize}
 \item \textbf{Analytical Stability and High-Fidelity Validation:} We present a formal stability analysis for the closed-loop dynamics, based on the system’s Lyapunov structure. The feasibility of the approach is validated through high-fidelity simulations that demonstrate oscillation attenuation and transport robustness, laying the theoretical groundwork for future experimental implementations in actual flight.
\end{itemize}

The remainder of this paper is organized as follows. Section~\ref{sec:methodology} presents the proposed methodology, including the geometric modeling and the hierarchical payload-oriented control architecture. Section~\ref{sec:simulations} reports the simulation and experimental results. Finally, Section~\ref{sec:conclusiones} concludes the paper and outlines future research directions.

%%%%%%%%%%%%%%%%%%%%%%%%%%%%%%%%%%%
\section{Methodology}\label{sec:methodology}
% REV4-C1: Agregar un párrafo introductorio de la Sección 2 que explique el flujo metodológico: modelo geométrico, control cinemático, compensador dinámico y análisis de estabilidad.

\subsection{Geometric Modeling of the UAV--Payload Relative Motion}
\label{sec:geom_model}

The geometric relationship between the UAV and the suspended payload is first introduced to formulate the outer-loop controller at the payload level. Figure \ref{fig:drone_carga} shows the payload transport system, with the defined variables and reference frames.
% REV4-C3: Unificar "Figure" con el resto de citas. Si se adopta "Fig.", cambiar aquí a "Fig.~\ref{fig:drone_carga}".
Let $\mathbf{x}\in\mathbb{R}^3$ denote the UAV position and $\mathbf{x}_l\in\mathbb{R}^3$ the payload position, both expressed in the inertial frame. The cable vector is defined as
\begin{equation}
\mathbf{h} = \mathbf{x} - \mathbf{x}_l,
\label{eq:cable_vector}
\end{equation}
where $l=\|\mathbf{h}\|$ is the cable length, assumed constant.

% REV1-C7: Mantener esta definición de $\mathbf{h}$ en todo el manuscrito. En la sección de métricas se usa $\mathbf{h}=p_L-p_U$, que tiene signo opuesto. Unificar la convención y ajustar fórmulas/captions si corresponde.

\begin{figure}
    \centering
    \includegraphics[width=0.5\linewidth]{Figures/drone_carga.eps}
    \caption{Drone-based payload transport system}
    \label{fig:drone_carga}
\end{figure}

For geometric interpretation, the relative payload position may also be described in a body-related frame attached to the UAV. Let ${}^{b}\mathbf{r}_l\in\mathbb{R}^3$ denote the vector from the UAV to the payload expressed in that frame. Then, the payload position in the inertial frame can be written as
\begin{equation}
\mathbf{x}_l = \mathbf{x} + \mathbf{R}_{\psi}\,{}^{b}\mathbf{r}_l,
\label{eq:payload_position_yaw}
\end{equation}
where $\mathbf{R}_{\psi}\in\mathbb{R}^{3\times 3}$ is the rotation matrix associated with the UAV yaw angle $\psi$, namely,
\begin{equation}
\mathbf{R}_{\psi}=
\begin{bmatrix}
\cos\psi & -\sin\psi & 0\\
\sin\psi & \cos\psi  & 0\\
0        & 0         & 1
\end{bmatrix}.
\label{eq:yaw_rotation}
\end{equation}
Differentiating \eqref{eq:payload_position_yaw} yields
\begin{equation}
\dot{\mathbf{x}}_l
=
\dot{\mathbf{x}}
+
\dot{\mathbf{R}}_{\psi}\,{}^{b}\mathbf{r}_l
+
\mathbf{R}_{\psi}\,{}^{b}\dot{\mathbf{r}}_l.
\label{eq:payload_velocity_yaw}
\end{equation}

From a symmetry perspective, the modularity of this architecture allows the quadrotor to maintain its structural symmetry (differential flatness properties) while the outer loop explicitly manages the asymmetric pendulum-like dynamics of the payload. This separation ensures that control actions intended for path-following do not conflict with the stabilization of the load.

Equations \eqref{eq:payload_position_yaw}--\eqref{eq:payload_velocity_yaw} show that the payload motion is geometrically influenced by the vehicle heading. This representation is useful to interpret the relative UAV--payload configuration, particularly in the horizontal plane. Nevertheless, the proposed swing-compensation mechanism does not explicitly use the yaw variable. The reason is that oscillation attenuation is primarily governed by horizontal translational motion and cable-state variation, whereas the yaw angle does not provide direct information relevant to the swing-reduction objective considered here.

Therefore, for control design and implementation purposes, the compact inertial-frame representation obtained directly from \eqref{eq:cable_vector} is adopted, namely,
\begin{equation}
\mathbf{x}_l = \mathbf{x} - \mathbf{h},
\qquad
\dot{\mathbf{x}}_l = \dot{\mathbf{x}} - \dot{\mathbf{h}}.
\label{eq:payload_velocity_compact}
\end{equation}
This form retains the essential geometric information required for payload-level trajectory tracking and swing attenuation while avoiding unnecessary coordinate transformations.

\subsection{Kinematic Control}\label{sec:kinematic}

This paper presents a hierarchical, payload-oriented control system for a drone. The study aims to track the desired trajectory and minimize oscillations of the suspended payload during the task. The proposed approach consists of two stages: a kinematic controller and a dynamic compensator. Figure \ref{fig:diagrama_controlador} shows the complete control structure diagram.

This section describes the kinematic controller that enables the payload to follow a trajectory and minimizes its oscillations. Before deriving the control laws, the available information for the controller, the operating assumptions, and the control objective are explicitly defined. The higher-level modules (mission management and trajectory planning) are not part of the controller synthesis; they only provide a feasible payload reference along with the associated operational constraints.

\begin{assumption}\label{as:physical}
The disturbance term $\boldsymbol{\delta}(t)$ acting on the velocity dynamics of the UAV and its time derivative are bounded, i.e., $\|\boldsymbol{\delta}(t)\|\le\delta_o$ and $\|\dot{\boldsymbol{\delta}}(t)\|\le\bar{\delta}$, and the normalized actuator commands operate within their linear (unsaturated) range during nominal operation. Physically, this assumption is justified by the finite energy of atmospheric disturbances. Wind gusts in urban delivery environments exhibit bounded magnitudes and limited rates of change, ensuring that the control gains and adaptation laws remain within feasible operational ranges.
\end{assumption}
\begin{figure}
    \centering
    \includegraphics[width=0.99\linewidth]{Figures/diagrama_controlador.eps}
    \caption{The kinematics layer computes payload-based reference velocities to reduce oscillations, whereas the dynamic controller guarantees strong tracking despite uncertainties in the system model and external disturbances. The whole design relies on feedback of the payload state and feedforward from the reference trajectory.}
    \label{fig:diagrama_controlador}
\end{figure}

It is important to note that the UAV's yaw angle is not specifically considered in the proposed payload control strategy. This is because the payload's dynamic behavior is primarily affected by the UAV's translational motion, while yaw has a negligible effect on the suspended payload's roll motion. Therefore, yaw is controlled separately using a direct reference tracking controller, which is decoupled from the proposed kinematic control strategies.
The outer-loop controller is formulated at the payload level so that the suspended load, rather than the UAV itself, tracks the prescribed reference trajectory. Let $\mathbf{x}_d\in\mathbb{R}^3$ denote the desired payload position, and let $\dot{\mathbf{x}}_d$ and $\ddot{\mathbf{x}}_d$ denote the corresponding reference velocity and acceleration. The payload tracking error is defined as
\begin{equation}
\tilde{\mathbf{x}} = \mathbf{x}_d - \mathbf{x}_l,
\label{eq:payload_tracking_error}
\end{equation}
with time derivative
\begin{equation}
\dot{\tilde{\mathbf{x}}} = \dot{\mathbf{x}}_d - \dot{\mathbf{x}}_l.
\label{eq:payload_tracking_error_dot}
\end{equation}
A nominal translational command is first introduced as
\begin{equation}
\mathbf{v}_{\mathrm{base}}
=
\dot{\mathbf{x}}_d
+
K_{\mathrm{traj}}\,\tilde{\mathbf{x}},
\label{eq:base_reference}
\end{equation}
where $K_{\mathrm{traj}}\in\mathbb{R}^{3\times 3}$ is a positive definite diagonal gain matrix. Although \eqref{eq:base_reference} ensures payload-level tracking, it does not compensate for the oscillations of the pendular configuration. Such compensation requires the cable vector $\mathbf{h}=\mathbf{x}-\mathbf{x}_l$ together with its rate $\dot{\mathbf{h}}$. Only $\mathbf{h}$ is measured; $\dot{\mathbf{h}}$ is reconstructed by a two-stage cable-rate estimator. A high-gain stage first produces a raw velocity estimate from the tracking error $\mathbf{h}-\hat{\mathbf{h}}$, 
\begin{equation}\label{eq:h_raw} \mathbf{h}_{p,\mathrm{raw}} = \frac{1}{\varepsilon_1}\,K_1\big(\mathbf{h}-\hat{\mathbf{h}}\big), 
\end{equation} 
with $K_1\in\mathbb{R}^{3\times 3}$ positive definite and $\varepsilon_1>0$. The internal state $\hat{\mathbf{h}}$ is then driven by the \emph{filtered} rate estimate $\hat{\mathbf{h}}_p$, 
\begin{equation}\label{eq:h_hat_dyn} \dot{\hat{\mathbf{h}}} = \hat{\mathbf{h}}_p, \end{equation} 
while $\hat{\mathbf{h}}_p$ is obtained by low-pass filtering the raw signal, 
\begin{equation}\label{eq:h_p_filter} \tau_h\,\dot{\hat{\mathbf{h}}}_p + \hat{\mathbf{h}}_p = \mathbf{h}_{p,\mathrm{raw}}, \qquad \tau_h>0. 
\end{equation} 
In this closed-loop structure, a smaller $\varepsilon_1$ yields faster tracking of the measurement (and thus a more responsive raw velocity estimate), whereas the time constant $\tau_h$ attenuates the noise amplified by differentiation. The filtered output $\hat{\mathbf{h}}_p$ is the estimate of $\dot{\mathbf{h}}$ used in the feedback law; at steady state $\hat{\mathbf{h}}\to\mathbf{h}$ and $\hat{\mathbf{h}}_p\to\dot{\mathbf{h}}$. The selection of $\varepsilon_1$ and $\tau_h$ therefore reflects a trade-off between estimation velocity and noise rejection. Since the dominant oscillations are excited by horizontal accelerations, only the horizontal cable component is compensated. Let $\mathbf{h}_h=[h_x,h_y]^T$ denote the horizontal cable displacement and $\hat{\mathbf{h}}_{p,h}=[\hat{h}_{p,x}, \hat{h}_{p,y}]^T$ its estimated rate. Under the convention $\mathbf{h}=\mathbf{x}-\mathbf{x}_l$, the quasi-static cable reference is \begin{equation}\label{eq:href} \mathbf{h}_h^\star = -\frac{l}{g}\,\ddot{\mathbf{x}}_{d,h}. \end{equation} This expression represents the \emph{equilibrium} inclination that the pendulum must adopt for the payload to follow the prescribed trajectory; it is not an oscillation to be suppressed. Both $\ddot{\mathbf{x}}_{d,h}$ and the reference rate $\dot{\mathbf{h}}_h^\star$ are obtained by filtered numerical differentiation of the (sufficiently smooth) reference trajectory, and $g$ is the gravitational acceleration. The swing-compensation term is then defined as \begin{equation}\label{eq:swing_compensation} \mathbf{v}_{\mathrm{comp}} = \begin{bmatrix} -K_h\big(\mathbf{h}_h-\mathbf{h}_h^\star\big) -K_{dh}\big(\hat{\mathbf{h}}_{p,h}-\dot{\mathbf{h}}_h^\star\big)\\[1mm] 0 \end{bmatrix}, \end{equation} with $K_h,K_{dh}\in\mathbb{R}^{2\times 2}$ positive definite and diagonal. The first term is a restoring action toward $\mathbf{h}_h^\star$: by referencing the equilibrium inclination rather than the pure vertical ($\mathbf{h}_h=\mathbf{0}$), it penalizes only the genuine swing (i.e., the deviation between the actual cable inclination and the inclination justified by the maneuver) instead of the inclination that the commanded acceleration itself requires. The second term damps the pendulum transients. A smooth start-up factor $\sigma(t)=1-e^{-t/T_s}$, with $T_s>0$, prevents impulsive commands at the onset of the maneuver, so that \begin{equation}\label{eq:final_reference} \mathbf{v}_c = \sigma(t)\,\big(\mathbf{v}_{\mathrm{base}} + \mathbf{v}_{\mathrm{comp}}\big). \end{equation} 
The yaw channel is treated separately, as it is not coupled to the dominant horizontal oscillations: 
\begin{equation}\label{eq:yaw_reference} \dot{\psi}_c = \dot{\psi}_d + k_{\psi}\tanh\!\big(k_{\psi}^t\,e_{\psi}\big), \qquad e_{\psi}=\psi_d-\psi, \end{equation} 
with $k_{\psi},k_{\psi}^t>0$. The complete outer-loop kinematic command is $\mathbf{u}_k=[\mathbf{v}_c^T,\ \dot{\psi}_c]^T$, which serves as the reference for the dynamic compensator presented in the next section.

The core philosophy of this symmetry-preserving approach is to transform the asymmetric, coupled dynamics into a predictable, decoupled tracking problem. The UAV acts as a symmetric force actuator, and the compensator's role is to neutralize the directional 'pull' of the payload and wind, maintaining a rotationally invariant tracking performance."
%==============================================================================
\subsection{Closed-Loop Error Analysis}
\label{sec:closed_loop_analysis}

The outer-loop error analysis is carried out under the perfect velocity-tracking
condition $\dot{\mathbf{x}}=\mathbf{u}_k$, so that the kinematic command is realized exactly by the inner loop. The analysis is supported by the following assumptions, which formalize the operating regime already invoked informally in the controller design.

\begin{assumption}[Fast cable-rate estimation]\label{as:estimation}
The two-stage estimator \eqref{eq:h_raw}--\eqref{eq:h_p_filter} operates on a fast
time scale ($\varepsilon_1,\tau_h$ small), so that the estimation residual
$\boldsymbol{\xi}_e=\hat{\mathbf{h}}_p-\dot{\mathbf{h}}$ is bounded and decays to a small neighborhood of the origin after a fast transient. This assumption is justified by the separation of time-scales between the high-frequency estimation loop and the slower mechanical response of the UAV. By selecting sufficiently small filter time constants ($(\epsilon_1, \tau_h)$), the estimation residual becomes negligible relative to the control sampling period, enabling a decoupled analysis of the task-space tracking error
\end{assumption}

\begin{assumption}[Quasi-static small inclination]\label{as:smallangle}
The horizontal cable displacement satisfies $\|\mathbf{h}_h\|/l\ll 1$, so that the
equilibrium reference $\mathbf{h}_h^\star=-(l/g)\,\ddot{\mathbf{x}}_{d,h}$ is exact to first order and the higher-order pendular terms are negligible. This assumption is consistent with the operational requirements of last-mile delivery, where cargo integrity is prioritized. By avoiding aggressive maneuvers, the system remains within a near-equilibrium regime that preserves the validity of the first-order pendular model while ensuring the safety of sensitive payloads during transport.
\end{assumption}

\begin{assumption}[Smooth bounded reference]\label{as:reference}
The reference trajectory is feasible and sufficiently smooth, with
$\ddot{\mathbf{x}}_d,\dddot{\mathbf{x}}_d\in\mathcal{L}_\infty$. Consequently
$\mathbf{h}_h^\star$ and $\dot{\mathbf{h}}_h^\star$ are well defined and bounded. Physically, this bounded smoothness requirement ensures that the reference motion does not demand infinite motor torque or instantaneous thrust changes. It reflects real-world flight planning where jerk and acceleration are limited to satisfy actuator constraints and avoid exciting the higher-order structural modes of the flexible cable.
\end{assumption}

\begin{assumption}[Vertical decoupling]\label{as:vertical}
The vertical cable motion $\dot{h_z}$ is regulated by the inner loop and remains
bounded; hence, the swing-compensation analysis is restricted to the horizontal
subspace $\tilde{\mathbf{h}}_h\in\mathbb{R}^2$. Vertical decoupling is justified by the quadrotor’s high thrust-to-weight ratio along the vertical axis ($(z_b)$). Since the vehicle possesses direct actuation to compensate for gravity and vertical payload reactions, the translational swing dynamics are primarily governed by horizontal accelerations, justifying a focused analysis on the horizontal swing-reduction subspace.
\end{assumption}

\noindent
These assumptions collectively ensure that the tracking problem remains symmetrical across the translational axes. By treating external perturbations as bounded variations around a stable operating point, the proposed architecture can effectively restore the system's nominal performance symmetry even in the presence of unmodeled dynamic effects.


Under perfect velocity tracking, substituting \eqref{eq:final_reference} into
\eqref{eq:payload_velocity_compact} and neglecting the start-up factor after the
initial transient ($\sigma(t)\to 1$) yields, for the horizontal cable error
$\tilde{\mathbf{h}}_h=\mathbf{h}_h-\mathbf{h}_h^\star$ and the payload error
$\tilde{\mathbf{x}}=\mathbf{x}_d-\mathbf{x}_l$, the \emph{first-order} cascade
\begin{align}
\dot{\tilde{\mathbf{h}}}_h
&= -K_h\,\tilde{\mathbf{h}}_h + \boldsymbol{\xi}_h,
\label{eq:cl_cable}\\[1mm]
\dot{\tilde{\mathbf{x}}}
&= -K_{\mathrm{traj}}\,\tilde{\mathbf{x}}
   + B\,\tilde{\mathbf{h}}_h
   + \boldsymbol{\zeta},
\label{eq:cl_payload}
\end{align}
where $B=\operatorname{diag}(K_h,0)\in\mathbb{R}^{3\times2}$ incorporates the horizontal action in the payload channel, $\boldsymbol{\xi}_h$ collects the bounded estimation residue of the assumption~\ref{as:estimation} together with the reference rate mismatch $\dot{\mathbf{h}}_h-\dot{\mathbf{h}}_h^\star$, and $\boldsymbol{\zeta}$ includes the bounded vertical contribution of the assumption~\ref{as:vertical} and the start-up residual term. The damping gain $K_{dh}$ does not appear in \eqref{eq:cl_cable} because it acts as a transient shaping practical term and is not required for stability, which is set later with $K_{dh}=0$.

The structure \eqref{eq:cl_cable}--\eqref{eq:cl_payload} is a cascade of two
input-to-state stable (ISS) first-order subsystems, and is summarized in the following
result.

\begin{theorem}[Outer-loop stability]\label{th:outer}
Consider the closed-loop error dynamics \eqref{eq:cl_cable}--\eqref{eq:cl_payload}
under Assumptions~\ref{as:estimation}--\ref{as:vertical}, with $K_h,K_{\mathrm{traj}}$
symmetric positive definite. Then:
\begin{enumerate}
  \item[\textnormal{(a)}] If $\boldsymbol{\xi}_h=\mathbf{0}$ and $\boldsymbol{\zeta}=\mathbf{0}$,
  the origin $(\tilde{\mathbf{x}},\tilde{\mathbf{h}}_h)=(\mathbf{0},\mathbf{0})$ is
  globally exponentially stable provided the gains satisfy
  \begin{equation}
  \lambda_{\min}(K_{\mathrm{traj}})>\tfrac{1}{2},\qquad
  \mu\,\lambda_{\min}(K_h)>\tfrac{1}{2}\|B\|^2,
  \label{eq:gain_condition}
  \end{equation}
  for some weight $\mu>0$.
  \item[\textnormal{(b)}] If $\boldsymbol{\xi}_h$ and $\boldsymbol{\zeta}$ are bounded,
  the cascade is ISS: all error signals remain bounded and converge to a compact
  neighborhood of the origin whose size is proportional to
  $\sup_t\|(\boldsymbol{\xi}_h,\boldsymbol{\zeta})\|$.
\end{enumerate}
\end{theorem}

\begin{proof}
Consider the weighted Lyapunov function
\begin{equation}
V=\tfrac{1}{2}\,\tilde{\mathbf{x}}^\top\tilde{\mathbf{x}}
 +\tfrac{\mu}{2}\,\tilde{\mathbf{h}}_h^\top\tilde{\mathbf{h}}_h,
 \qquad \mu>0,
\label{eq:cl_lyapunov}
\end{equation}
which is positive definite and radially unbounded. Its time derivative along
\eqref{eq:cl_cable}--\eqref{eq:cl_payload} is
\begin{equation}
\dot V=
-\tilde{\mathbf{x}}^\top K_{\mathrm{traj}}\tilde{\mathbf{x}}
-\mu\,\tilde{\mathbf{h}}_h^\top K_h\tilde{\mathbf{h}}_h
+\tilde{\mathbf{x}}^\top B\,\tilde{\mathbf{h}}_h
+\tilde{\mathbf{x}}^\top\boldsymbol{\zeta}
+\mu\,\tilde{\mathbf{h}}_h^\top\boldsymbol{\xi}_h.
\label{eq:cl_Vdot}
\end{equation}

\emph{Part (a).} With $\boldsymbol{\xi}_h=\boldsymbol{\zeta}=\mathbf{0}$, the cross term
is bounded by Young's inequality,
\begin{equation}
\tilde{\mathbf{x}}^\top B\,\tilde{\mathbf{h}}_h
\le \tfrac{1}{2}\|\tilde{\mathbf{x}}\|^2
   +\tfrac{1}{2}\|B\|^2\|\tilde{\mathbf{h}}_h\|^2,
\label{eq:young}
\end{equation}
so that
\begin{equation}
\dot V \le
-\Big(\lambda_{\min}(K_{\mathrm{traj}})-\tfrac{1}{2}\Big)\|\tilde{\mathbf{x}}\|^2
-\Big(\mu\,\lambda_{\min}(K_h)-\tfrac{1}{2}\|B\|^2\Big)\|\tilde{\mathbf{h}}_h\|^2.
\label{eq:cl_Vdot_bound}
\end{equation}
Under \eqref{eq:gain_condition} both coefficients are strictly positive, hence
$\dot V\le -\alpha V$ for some $\alpha>0$, which proves global exponential stability of
the origin.

\emph{Part (b).} With bounded $\boldsymbol{\xi}_h,\boldsymbol{\zeta}$, applying Young's
inequality to the disturbance terms as well gives
\begin{equation}
\dot V \le
-c_1\|\tilde{\mathbf{x}}\|^2
-c_2\|\tilde{\mathbf{h}}_h\|^2
+\rho\big(\|\boldsymbol{\xi}_h\|^2+\|\boldsymbol{\zeta}\|^2\big),
\label{eq:cl_Vdot_iss}
\end{equation}
with $c_1,c_2,\rho>0$ obtained from \eqref{eq:gain_condition} with a margin reserved
for the disturbance bounds. Thus $\dot V<0$ outside a compact set, which establishes
input-to-state stability of the cascade with respect to
$(\boldsymbol{\xi}_h,\boldsymbol{\zeta})$. Equivalently, since subsystem
\eqref{eq:cl_cable} is ISS with respect to $\boldsymbol{\xi}_h$ and subsystem
\eqref{eq:cl_payload} is ISS with respect to $(\tilde{\mathbf{h}}_h,\boldsymbol{\zeta})$, the cascade interconnection of ISS systems is ISS, recovering the same conclusion.
This proves part~(b) and completes the proof.
\end{proof}

\noindent
Theorem~\ref{th:outer} makes explicit the design trade-off hidden in the original
formulation: the trajectory gain $K_{\mathrm{traj}}$ sets the payload-tracking rate,
while the swing gain $K_h$ must be large enough, relative to the coupling $\|B\|$ and
the chosen weight $\mu$, to dominate the restoring action injected into the payload
channel. The damping gain $K_{dh}$, although not required for stability, is retained in
the implementation to shape the pendular transient and reduce overshoot, as confirmed
by the simulation results.

\subsection{Dynamic Compensation}\label{sec:dynamic}

This section presents a robust adaptive dynamic compensator that enforces the perfect velocity-tracking condition $\dot{\mathbf{x}}=\mathbf{u}_k$ assumed in the previous section, driving the velocity-tracking error to a small neighborhood of the origin despite model uncertainty and external disturbances. A Super-Twisting algorithm is adopted because, as a second-order sliding-mode technique, it preserves the robustness and finite-time convergence of sliding-mode control while substantially attenuating the chattering associated with first-order sliding-mode controllers. In addition, an adaptation law for the model parameters is introduced to reduce the conservativeness of the control gains and improve tracking performance under varying payload conditions.

As a first step, the dynamic model of the vehicle is presented. The one proposed in \cite{Lucas_ICUAS2014}, formulated for a commercial drone with internal stabilization loops, was adopted, since it directly provides the mapping between the normalized input commands and the vehicle accelerations, a relationship that constitutes the basis for the design of the dynamic compensator.
\begin{equation}\label{eq:ModeloLucas}
    \ddot{\mathbf{x}} = \boldsymbol{\kappa}_u\mathbf{u} - \boldsymbol{\kappa}_v\dot{\mathbf{x}} +\boldsymbol{\delta}(t)
\end{equation}

% REV1-C7/C8: Revisar el signo de $\boldsymbol{\delta}(t)$. En ecuaciones posteriores aparece con signo positivo/negativo de forma no consistente. Además, describir físicamente qué incluye la perturbación y qué perfiles/magnitudes se usan en simulación.
\noindent where $\boldsymbol{\kappa}_u \in \mathbb{R}^{n\times n}$ and $\boldsymbol{\kappa}_v \in \mathbb{R}^{n\times n}$, being \textit{n=4}. The state vector $\mathbf{x}\in \mathbb{R}^{n\times 1}$, and the input vector $\mathbf{u}\in \mathbb{R}^{n\times 1}$ with inputs $[-1,+1]$ that represent the normalized commands sent to the vehicle, these can be velocities or reference angles. In addition, the second term considers all aerodynamic effects that affect the vehicle and that are related to the vehicle velocity. Finally, external disturbances and modeling errors are represented in the $\boldsymbol{\delta}(t)$ term.
By hypothesis, it is considered that the derivatives of this perturbation are bounded. Therefore, $\| \boldsymbol{\delta}(t)\| \leq \boldsymbol{\delta}_o$. The kinematic controller commands are reference velocities $v_d$. Therefore, the following sliding surface is proposed:
% REV3-C13/C14: Describir explícitamente las perturbaciones externas variantes en el tiempo usadas en simulación y distinguir casos periódicos y no periódicos. Esta hipótesis debe conectarse con los perfiles simulados.

\begin{equation}
    \mathbf{s} = \Tilde{\mathbf{v}} + \boldsymbol{\gamma}\int \Tilde{\mathbf{v}} dt
\end{equation}
Deriving the above equation,

\begin{equation}
    \dot{\mathbf{s}} = \dot{\Tilde{\mathbf{v}}} + \boldsymbol{\gamma} \Tilde{\mathbf{v}} = (\ddot{\mathbf{x}} - \dot{\mathbf{v}}_d) + \boldsymbol{\gamma}(\dot{\mathbf{x}} - \mathbf{v}_d)
\end{equation}

The equivalent control law in the sliding mode can be found by solving $\dot{\mathbf{s}}=0$. Therefore, \eqref{eq:ModeloLucas} is replaced in the above equation,

\begin{equation}\label{eq:Condicion_S}
    \dot{\mathbf{s}} = \big(\boldsymbol{\kappa}_u\mathbf{u} - \boldsymbol{\kappa}_v\dot{\mathbf{x}}+\boldsymbol{\delta}(t)\big) - 
    \dot{\mathbf{v}}_d - \boldsymbol{\gamma} \Tilde{\mathbf{v}} = 0
\end{equation}

clearing the control action $\mathbf{u}$, it can be expressed as:

\begin{eqnarray}\label{eq:AccionIdeal}
    \mathbf{u} &=& \boldsymbol{\kappa}_u^{-1}( -\boldsymbol{\gamma} \tilde{\mathbf{v}} + \dot{\mathbf{v}}_d + \boldsymbol{\kappa}_v v) +\boldsymbol{\kappa}_u^{-1}\boldsymbol{\delta}(t) 
\end{eqnarray}

The above equation represents the ideal control action, but the values of $\boldsymbol{\delta}(t)$ and the actual values of the $\boldsymbol{\kappa}_u$ and $\boldsymbol{\kappa}_v$ matrices are unknown. Therefore, the following super-twisting type control is proposed for the previously presented control objective

\begin{eqnarray}\label{eq:AccionControl}
    \mathbf{u} &=& \hat{\boldsymbol{\kappa}}_u^{-1}( -\boldsymbol{\gamma} \Tilde{\mathbf{v}} + \dot{\mathbf{v}}_d + \hat{\boldsymbol{\kappa}}_v v) - \mathbf{w} = \boldsymbol{\phi \hat{\theta}} - \mathbf{w}\nonumber\\
    \mathbf{w} &=& \mathbf{K}_1\sqrt{||\mathbf{s}||}\sign(\mathbf{s}) + \int \mathbf{K}_2\sign(\mathbf{s}) dt
\end{eqnarray}

The matrices $\hat{\boldsymbol{\kappa}}_u$ and $\hat{\boldsymbol{\kappa}}_v$ represent the matrices with known parameters, and the $\boldsymbol{\phi \hat{\theta}}$ term is the parameterization of part of the control law. The estimated parameters are denoted by $\boldsymbol{\hat{\theta}}$.
The control law \eqref{eq:AccionControl} is linearly parameterized in the unknown
model parameters. For the decoupled model of \cite{Lucas_ICUAS2014}, with
$\boldsymbol{\kappa}_u$ and $\boldsymbol{\kappa}_v$ diagonal, the parameterized term
$\boldsymbol{\phi}\boldsymbol{\theta}=\boldsymbol{\kappa}_u^{-1}\mathbf{a}
+\boldsymbol{\kappa}_u^{-1}\boldsymbol{\kappa}_v\dot{\mathbf{x}}$ admits, on each
channel $i\in\{1,\dots,4\}$, the two unknown parameters
\begin{equation}\label{eq:param_def}
\theta_{1,i}=\frac{1}{\kappa_{u,i}},\qquad
\theta_{2,i}=\frac{\kappa_{v,i}}{\kappa_{u,i}},
\end{equation}
so that $(\boldsymbol{\phi}\boldsymbol{\theta})_i=\theta_{1,i}\,a_i+\theta_{2,i}\,\dot{x}_i$,
with $\mathbf{a}=\dot{\mathbf{v}}_d+\boldsymbol{\gamma}\tilde{\mathbf{v}}$ a known signal.
Stacking the four channels, the parameter vector is
$\boldsymbol{\theta}\in\mathbb{R}^{p}$ with $p=8$, and the regressor
$\boldsymbol{\phi}\in\mathbb{R}^{4\times 8}$ has the block-diagonal structure
\begin{equation}\label{eq:regressor}
\boldsymbol{\phi}=
\begin{bmatrix}
a_1 & \dot{x}_1 & & & & & & \\
 & & a_2 & \dot{x}_2 & & & & \\
 & & & & a_3 & \dot{x}_3 & & \\
 & & & & & & a_4 & \dot{x}_4
\end{bmatrix},
\end{equation}
where empty entries are zero. Consequently $\mathbf{A}(t)=\boldsymbol{\Gamma}\,
\boldsymbol{\phi}(t)^T\boldsymbol{\phi}(t)\in\mathbb{R}^{8\times 8}$ and the adaptation
gain $\boldsymbol{\Gamma}\in\mathbb{R}^{8\times 8}$. Therefore, replacing \eqref{eq:AccionControl} in \eqref{eq:Condicion_S} gives,

\begin{eqnarray}
     \nonumber\dot{\mathbf{s}} &=& \boldsymbol{\kappa}_u \left(\boldsymbol{\phi \hat{\theta}} - \mathbf{K}_1\sqrt{||\mathbf{s}||}\sign(\mathbf{s}) + \int \mathbf{K}_2\sign(\mathbf{s}) dt \right) -...\\
    &-& \boldsymbol{\kappa}_v\dot{\mathbf{x}} - \boldsymbol{\delta}(t)  - 
    \dot{\mathbf{v}}_d + \boldsymbol{\gamma} \Tilde{\mathbf{v}}
\end{eqnarray}

We define $\mathbf{a} = \dot{\mathbf{v}}_d + \boldsymbol{\gamma} \Tilde{\mathbf{v}}$ to improve the readability of the mathematical development. Replacing this and operating on the above equation yields.

\begin{eqnarray}
    \dot{\mathbf{s}} &=& \boldsymbol{\kappa}_u \left(\boldsymbol{\phi \hat{\theta}} - \underbrace{\boldsymbol{\kappa}_u^{-1}\mathbf{a} - \boldsymbol{\kappa}_u^{-1}\boldsymbol{\kappa}_v\dot{\mathbf{x}}}_{\boldsymbol{\phi \theta}} \right)+...\\  
      \nonumber&+&\boldsymbol{\kappa}_u\left( \mathbf{K}_1\sqrt{||\mathbf{s}||}\sign(\mathbf{s}) + \int \mathbf{K}_2\sign(\mathbf{s}) dt  \right) - \boldsymbol{\delta}(t) 
\end{eqnarray}
The above equation results,

\begin{equation}\label{eq:SuperficieFinal}
    \dot{\mathbf{s}} = \boldsymbol{\kappa}_u \boldsymbol{\phi \Tilde{\theta}}  
     +\boldsymbol{\kappa}_u\left( \mathbf{K}_1\sqrt{||\mathbf{s}||}\sign(\mathbf{s}) + \int \mathbf{K}_2\sign(\mathbf{s}) dt  \right) - \boldsymbol{\delta}(t) 
\end{equation}

\begin{theorem}[Inner-loop robust stability]\label{th:inner}
Consider the velocity-error sliding dynamics \eqref{eq:SuperficieFinal} under the composite adaptive super-twisting control law \eqref{eq:AccionControl} and the adaptation law \eqref{eq:ajuste2}. Suppose that Assumption~\ref{as:physical} holds (so that $\|\dot{\boldsymbol{\delta}}(t)\|\le\bar{\delta}$) and that the super-twisting gain satisfies the matching condition $\mathbf{K}_2\ge\|\boldsymbol{\kappa}_u^{-1}\dot{\boldsymbol{\delta}}(t)\|$. Then the sliding variable $\mathbf{s}$ and the parameter error $\tilde{\boldsymbol{\theta}}$ are uniformly bounded, $\dot{V}\le 0$, and the velocity-tracking error converges to zero, $\tilde{\mathbf{v}}\to\mathbf{0}$. If, in addition, the regressor $\boldsymbol{\phi}(t)$ is persistently exciting, then $\tilde{\boldsymbol{\theta}}\to\mathbf{0}$.
\end{theorem}

\begin{proof}
The Lyapunov theory is used to perform the stability analysis. First, the following Lyapunov candidate function is proposed,

\begin{equation*}\label{eq:CandidataLyapunov}
    V = \cfrac{1}{2}\mathbf{s}^T\boldsymbol{\kappa}^{-1}\mathbf{s} + \cfrac{1}{2}\boldsymbol{\Tilde{\theta}}^T\boldsymbol{\gamma} \boldsymbol{\Tilde{\theta}} 
\end{equation*}

Deriving the above equation, it is obtained,

\begin{equation}
    \dot{V} = \mathbf{s}^T\boldsymbol{\kappa}^{-1}\dot{\mathbf{s}} + \boldsymbol{\Tilde{\theta}}^T\boldsymbol{\gamma} \boldsymbol{\dot{\Tilde{\theta}}} 
\end{equation}

Replacing \eqref{eq:SuperficieFinal} in the above equation

\begin{eqnarray}
   \nonumber \dot{V} &=& \mathbf{s}^T\boldsymbol{\kappa}^{-1} \bigg( \boldsymbol{\kappa}_u \boldsymbol{\phi \Tilde{\theta}}  
     +\boldsymbol{\kappa}_u\Big( \mathbf{K}_1\sqrt{||\mathbf{s}||}\sign(\mathbf{s}) +...\\
     &+& \int \mathbf{K}_2\sign(\mathbf{s}) dt \Big) - \boldsymbol{\delta}(t) \bigg) + \boldsymbol{\Tilde{\theta}}^T\boldsymbol{\gamma} \boldsymbol{\dot{\Tilde{\theta}}} \nonumber\\
    \nonumber \dot{V} &=& \mathbf{s}^T \boldsymbol{\phi \Tilde{\theta}} - \mathbf{s}^T \bigg(\mathbf{K}_1\sqrt{||\mathbf{s}||}\sign(\mathbf{s}) -...\\ 
     &-&\int \mathbf{K}_2\sign(\mathbf{s}) dt
     -\boldsymbol{\kappa}^{-1}_u\boldsymbol{\delta}(t)  \bigg)+ \boldsymbol{\Tilde{\theta}}^T\boldsymbol{\gamma} \boldsymbol{\dot{\Tilde{\theta}}}
\end{eqnarray}

By the property of transpose, we obtain the following equalities:

\begin{equation}
 \mathbf{s}^T \boldsymbol{\phi \Tilde{\theta}} = (\mathbf{s}^T \boldsymbol{\phi \Tilde{\theta}})^T =\boldsymbol{\Tilde{\theta}}^T \boldsymbol{\phi}^T\mathbf{s}
\end{equation}

Based on the properties of transpose and operand, it is obtained,

\begin{eqnarray}\label{eq:Lyapunto}
    \nonumber\dot{V} &=& \boldsymbol{\tilde{\theta}}^T \boldsymbol{\phi}^T\mathbf{s} + \boldsymbol{\tilde{\theta}}^T\boldsymbol{\gamma} \boldsymbol{\dot{\tilde{\theta}}} -...\\
    \nonumber&-&\mathbf{s}^T\left(\mathbf{K}_1\sqrt{||\mathbf{s}||}\sign(\mathbf{s}) - \int \mathbf{K}_2\sign(\mathbf{s}) dt
     -\boldsymbol{\kappa}^{-1}_u\boldsymbol{\delta}(t) \right) \nonumber\\
    \dot{V} &=& \boldsymbol{\tilde{\theta}}^T(\boldsymbol{\phi}^T\mathbf{s} + \boldsymbol{\gamma} \boldsymbol{\dot{\tilde{\theta}}}) -...\\
     \nonumber&-& \mathbf{s}^T\left(\mathbf{K}_1\sqrt{||\mathbf{s}||}\sign(\mathbf{s}) - \int \mathbf{K}_2\sign(\mathbf{s}) dt
     -\boldsymbol{\kappa}^{-1}_u\boldsymbol{\delta}(t) \right)
\end{eqnarray}

\noindent The following law of adaptation of the system parameters of the form is proposed:
\begin{equation}\label{eq:ajuste}
    \boldsymbol{\dot{\tilde{\theta}}} = -\boldsymbol{\gamma}^{-1}\boldsymbol{\phi}^T\mathbf{s}^T-\boldsymbol{\gamma}\boldsymbol{\phi}^T\boldsymbol{\phi}\boldsymbol{\tilde{\theta}}
\end{equation}
\noindent Defining $\mathbf{A}(t)=\boldsymbol{\gamma}\boldsymbol{\phi}(t)^T\boldsymbol{\phi}(t)$ as $\mathbf{A}(t)\geq0 \in \mathbb{R}^{8 \times 8},\forall t>0$. The matrix $\mathbf{A}$ is a symmetric matrix with diagonal elements $a_{i,j}\geq0$. Therefore, it is a positive semi-definite matrix and its eigenvalues are $\lambda_{i,j}\geq0, \forall i=j$, then eq. \eqref{eq:ajuste} can be expressed as:
\begin{equation}\label{eq:ajuste2}
   \boldsymbol{\dot{\tilde{\theta}}} = -\boldsymbol{\gamma}^{-1}\boldsymbol{\phi}^T\mathbf{s}^T-\mathbf{A}\boldsymbol{\tilde{\theta}}
\end{equation}
\noindent Replacing the adjustment law \eqref{eq:ajuste2} in eq. \eqref{eq:Lyapunto} leads to:
\begin{eqnarray}
   \nonumber \dot{V} &=& -\boldsymbol{\tilde{\theta}}^T\mathbf{A}\boldsymbol{\tilde{\theta}}-\mathbf{s}^T\mathbf{K}_1\sqrt{||\mathbf{s}||}\sign(\mathbf{s})-...\\ &-& \mathbf{s}^T\int \mathbf{K}_2\sign(\mathbf{s}) dt
     -\mathbf{s}^T\boldsymbol{\kappa}^{-1}_u\boldsymbol{\delta}(t)
\end{eqnarray}
Applying the triangular inequality to the above equation and taking into account we obtain,
\begin{eqnarray}
 \nonumber    \dot{V} &\leq& -\lambda_{min}(\mathbf{A})\|\boldsymbol{\tilde{\theta}}\|^2-\|\mathbf{K}_1\|\|\mathbf{s}^T\|\sqrt{\|\mathbf{s}\|}- ...\\
    &-& \|\mathbf{s}^T\|\left\|\int \mathbf{K}_2 dt\right\|
     -||\mathbf{s}^T|| \left\| \int \boldsymbol{\kappa}^{-1}_u \dot{\boldsymbol{\delta}}(t) dt \right\| \nonumber \\
 \nonumber  \dot{V} &\leq&-\lambda_{min}(\mathbf{A})\|\boldsymbol{\tilde{\theta}}\|^2 -\|\mathbf{K}_1\|\|\mathbf{s}^T\|\sqrt{\|\mathbf{s}\|} -...\\
  &-&  \|\mathbf{s}^T\| \left( \left\|\int \mathbf{K}_2 dt\right\|
     - \left\| \int \boldsymbol{\kappa}^{-1}_u \dot{\boldsymbol{\delta}}(t) dt \right\|\right)     
\end{eqnarray}
The stability condition results from the fact that $\mathbf{K}_2\geq \|\boldsymbol{\kappa}^{-1}_u \dot{\boldsymbol{\delta}}(t)\|$. By hypothesis, it was assumed that the velocities of the perturbation $\boldsymbol{\delta}(t)$ are bounded. Finally, it is obtained that
\begin{equation}\label{eq:DemoEstabilidad}
    \dot{V} \leq -\lambda_{min}(\mathbf{A})\|\boldsymbol{\tilde{\theta}}\|^2 -\|\mathbf{K}_1\|\|\mathbf{s}^T\|\sqrt{\|\mathbf{s}\|} \leq 0
\end{equation}
Since $\mathbf{A}(t)=\boldsymbol{\gamma}\boldsymbol{\phi}^T\boldsymbol{\phi}\geq\mathbf{0}$ and $\mathbf{K}_1 >\mathbf{0}$, inequality \eqref{eq:DemoEstabilidad} yields $\dot{V}\le 0$, which proves that $\mathbf{s}$ and $\tilde{\boldsymbol{\theta}}$ are uniformly bounded and that $V$ is non-increasing. The negative-definite dependence of $\dot{V}$ on $\mathbf{s}$ (through the term $-\|\mathbf{K}_1\|\,\|\mathbf{s}^T\|\sqrt{\|\mathbf{s}\|}$) together with Barbalat's lemma guarantees that the sliding variable, and hence the velocity-tracking error $\tilde{\mathbf{v}}$, converge to zero as $t\to\infty$. Because $\mathbf{A}(t)$ is only positive semidefinite, the parameter error $\tilde{\boldsymbol{\theta}}$ is guaranteed to remain bounded; its convergence to zero additionally requires the regressor $\boldsymbol{\phi}(t)$ to be persistently exciting. 
\end{proof}



\section{Simulation Results}\label{sec:simulations}

The proposed control system was validated using a high-fidelity simulator for cable-suspended cargo transport presented in \cite{salinas2023unified}. The quadrotor dynamics are implemented in C++ with the parameters of \cite{michael2010grasp} %achtelik2010nonlinear}
, the cable is discretized as a chain of mass-spring-damper elements, and the package is modeled as a cubic body subject to aerodynamic drag. Table ~\ref{tab:sim_params} summarizes the main physical and simulation settings common to all the described scenarios.

\begin{table}[ht]
\centering
\caption{Main simulation parameters.}
\label{tab:sim_params}
\renewcommand{\arraystretch}{1.15}
\setlength{\tabcolsep}{5pt}
\begin{tabular}{l c l}
\toprule
\textbf{Parameter} & \textbf{Symbol} & \textbf{Value} \\
\midrule
\multicolumn{3}{l}{\textit{Quadrotor UAV} \cite{michael2010grasp}} \\ %achtelik2010nonlinear}} \\
UAV mass & $m$ & $0.48\,\mathrm{kg}$ \\
Principal moments of inertia & $[I_{xx}; I_{yy}; I_{zz}]$ & $[5.6;\ 5.6;\ 8.1]\times10^{-3}\,\mathrm{kg{\cdot}m^2}$ \\
Rotor-arm length & $L$ & $0.34\,\mathrm{m}$ \\
Rotor thrust coefficient & $k_F$ & $6.11\times10^{-8}\,\mathrm{N/(r/min)^2}$ \\
\midrule
\multicolumn{3}{l}{\textit{Cable and payload}} \\
Payload mass & $m_L$ & $0.10\,\mathrm{kg}$ \\
Payload size (cube edge) & $d_L$ & $0.10\,\mathrm{m}$ \\
Cable nominal length & $l$ & $2.0\,\mathrm{m}$ \\
Cable segments & $N_c$ & $40$ \\
Link mass & $m_{\mathrm{link}}$ & $0.001\,\mathrm{kg}$ \\
Spring stiffness & $k_c$ & $5000\,\mathrm{N/m}$ \\
Spring damping & $c_c$ & $0.6\,\mathrm{N{\cdot}s/m}$ \\
Payload drag coefficient & $C_d$ & $1.05$ \\
\midrule
\multicolumn{3}{l}{\textit{Simulation settings}} \\
Control sampling period & $T_c$ & $0.02\,\mathrm{s}$ ($50\,\mathrm{Hz}$) \\
Gravitational acceleration & $g$ & $9.81\,\mathrm{m/s^2}$ \\
Air density & $\rho$ & $1.151\,\mathrm{kg/m^3}$ \\
Normalized actuator commands & $\mathbf{u}$ & $[-1,\,+1]$ \\
\bottomrule
\end{tabular}

\vspace{0.15cm}
%\footnotesize{The payload is modeled as a cubic body with aerodynamic drag; the cable is discretized as a mass--spring--damper chain, so distributed flexibility and tension transients enter the simulator as unmodeled dynamics relative to the quasi-static control model.}
\end{table}

The control objective is for the load to follow a reference trajectory in space. For controller validation, no information from the dynamic model was considered, nor was any identification performed. The controller parameters were initialized with arbitrary values of the nominal matrices $\hat{\boldsymbol{\kappa}}_u = [0.5, 0.5, 0.5, 0.5]^T$ and $\hat{\boldsymbol{\kappa}}_v = [0.75, 0.75, 0.75, 0.75]^T$. Therefore, even in nominal scenarios, the closed loop is exercised under parametric mismatch and unmodeled dynamics because the actual vehicle parameters differ from their nominal counterparts, and the simulator preserves the distributed cable flexibility, tension transients, and pendulum coupling beyond the quasi-static model of the assumption ~\ref{as:physical}.

Two simulation scenarios were defined. In the first group, an ascending circular reference trajectory was used to compare the system's behavior with only the kinematic controller and with the dynamic compensator. In the second group, a more demanding reference was considered, and prescribed wind disturbances were applied over two time intervals. In this test, the aerodynamic disturbances acted on the UAV, the cable, and the payload. The disturbance on the payload was calculated from its physical dimensions using the cubic body drag model, rather than a point mass approximation. All runs were deterministic, with fixed references and prescribed wind profiles.

To evaluate the performance of the proposed control architecture, three complementary metrics are used to quantify the oscillatory behavior of the cable–payload system. From \eqref{eq:cable_vector}, the first metric is the cable inclination angle, defined as
% REV1-C7: Esta definición de $\mathbf{h}$ contradice \eqref{eq:cable_vector}. Unificar como $\mathbf{h}=\mathbf{x}-\mathbf{x}_l$ o renombrar el vector de métricas para evitar errores de signo.
\begin{equation}
\alpha(t)=\mathrm{atan2}\!\left(\sqrt{h_x^2+h_y^2},-h_z\right),
\end{equation}
which quantifies the extent to which the cable deviates from its vertical orientation and signifies the amplitude of the swinging payload. This criterion has been extensively used in various aerial cargo transport problems because it directly relates to the amplitude of the oscillation concerning the vertical position. The second metric is the horizontal displacement of the payload relative to the UAV, defined as
\begin{equation}
\rho(t)=\sqrt{h_x^2+h_y^2},
\end{equation}
which corresponds to the lateral distance of the payload from the vertical axis passing through the UAV. Finally, the dynamic behavior of the swing is characterized through the tangential relative velocity between the UAV and the payload. Let $\mathbf{v}_{rel}=\mathbf{v}_L-\mathbf{v}_U$ be the relative velocity and $\hat{\mathbf{h}}=\mathbf{h}/\|\mathbf{h}\|$ the unit cable direction. The tangential component is computed as
\begin{equation}
\mathbf{v}_{\perp}=\mathbf{v}_{rel}-(\mathbf{v}_{rel}\cdot\hat{\mathbf{h}})\hat{\mathbf{h}}.
\end{equation}
While $\alpha(t)$ and $\rho(t)$ describe the geometric magnitude of the oscillation, $\|\mathbf{v}_{\perp}\|$ captures the dynamic energy associated with the pendular motion. The joint analysis of these metrics provides a comprehensive assessment of both the amplitude and damping characteristics of the payload oscillations.

Tracking performance is evaluated using RMSE (Root Mean Squared Error) metrics, whereas oscillatory behavior is quantified using RMS (Root Mean Squared) values of the swing-related variables.

\subsection{Payload trajectory tracking}

A smooth, continuously differentiable reference trajectory was defined for the load to evaluate the system under representative air transport conditions without introducing discontinuities to the control signals.
\begin{eqnarray}
    \mathbf{x}_d = \begin{bmatrix} x_d \\ y_d \\ z_d \\ \psi_d\end{bmatrix}= \begin{bmatrix} 3 \sin{2\pi t}(m) \\  3 \cos{2\pi t} (m)\\ 6+\frac{t}{10} (m) \\ \arctan(\dot{y}_d / \dot{x}_d) \end{bmatrix}
\end{eqnarray}
It should be noted that in any trajectory-tracking problem in which there is a requirement for lateral acceleration, the rope cannot be kept vertical. Rather, there has to be some form of inclination. In the case of circular and/or elliptical paths, this is in the form of a nonzero quasi-static inclination that provides the necessary acceleration for the circular motion. Therefore, the raw quantities $\alpha(t)$ and $\rho(t)$ contain both a steady-state component imposed by the trajectory and a truly oscillatory component associated with payload swing. As a result, the direct use of $RMS_\alpha$ and $RMS_\rho$ may overestimate the actual oscillation level. This quasi-static component can be recognized as a common-mode effect in oscillation-related quantities, as it depends mainly on the reference trajectory rather than the control strategy. 

It was evaluated whether it was possible to achieve the load trajectory tracking task using only the kinematic controller \eqref{eq:base_reference}, which does not consider the oscillatory behavior of the suspended load. Multiple simulations were performed, confirming that this controller fails to meet the control objective and the system becomes unstable, as shown in Figure \ref{fig:CinematicoSolo_CinematicoCompensado}(a). Subsequently, the same simulation was performed using the controller proposed in \eqref{eq:final_reference} and \eqref{eq:yaw_reference}. In this case, the controller achieves trajectory tracking with very good control error and oscillation values, as shown in Figure \ref{fig:CinematicoSolo_CinematicoCompensado}(b). More specifically, the payload's movement becomes smoother, while the oscillations induced by the dynamics of the suspended load during task execution are significantly mitigated. This improvement in behavior can initially be observed in the time evolution of the system's main variables, as shown in Figure ~\ref{fig:CinematicoII_Variabes}. This suggests that the proposed kinematic compensation method is effective at counteracting the pendulum load perturbation on the vehicle.

\begin{figure}[t] \centering 
\setlength{\tabcolsep}{2pt} 
\begin{tabular}{@{}c@{\hspace{0.02\linewidth}}c@{\hspace{0.02\linewidth}}c@{}} \includegraphics[width=0.45\linewidth]{Figures/3D_Vista_Base.eps} & \includegraphics[width=0.45\linewidth]{Figures/3D_Vista_CinematicoII.eps} \\ \small\textbf{(a)} Unstable. & \small\textbf{(b)} Simulation completed successfully. 
\end{tabular}
\caption{In \textbf{(a)} it is observed that the basic pneumatic controller cannot stabilize the control system for reference trajectory tracking. The generated disturbances destabilize the drone-payload system. In \textbf{(b)}, it is observed that the proposed kinematic controller can stabilize the system and the trajectory tracking task is achieved.} \label{fig:CinematicoSolo_CinematicoCompensado} 
\end{figure}

The oscillation-related metrics shown in Figure ~\ref{fig:CinematicoII_Variabes} provide further evidence. Their evolution confirms a systematic decrease in all oscillation measurements throughout the maneuver. Taken together, these results demonstrate that the proposed kinematic oscillation reduction scheme is an effective alternative for improving payload transport performance.

\begin{figure}
\centering
\includegraphics[width=0.95\linewidth]{Figures/CinematicoII.eps}
\caption{Time evolution of the reference and actual payload states under the proposed kinematic swing-compensation scheme. From top to bottom, the plots show the position-tracking responses in the $x$, $y$, and $z$ directions, along with the yaw-angle evolution. The results indicate accurate trajectory tracking with a substantial reduction in the oscillatory effects induced by the suspended payload.}\label{fig:CinematicoII_Variabes}
\end{figure}

\begin{figure}
\centering
\includegraphics[width=0.95\linewidth]{Figures/CinematicoII_Angles.eps}
\caption{Time evolution of the oscillation-related performance metrics obtained with the proposed kinematic swing-compensation scheme: payload swing angle $\alpha$, payload position error magnitude $\rho$, and tangential velocity norm $\|v_\perp\|$. The results show that, after the initial transient, the oscillatory motion remains bounded and significantly attenuated throughout the maneuver.}\label{fig:CinematicoII_Metricas}
\end{figure}

\subsection{Dynamic Compensator Performance and Resulting Closed-Loop Behavior}
\label{subsec:dynamic_compensator_results}

The next step was to include the dynamic compensator according to the block diagram presented in Figure \ref{fig:diagrama_controlador} and analyze the behavior of the entire control loop proposed in this work. In this context, the function of the dynamic compensation layer is not to directly suppress payload oscillations, but rather to ensure that the vehicle accurately follows the reference signals generated by the external loop kinematic controller. This is especially relevant in suspended load transport, as any discrepancy between the commanded and actual movement of the UAV can cause additional excitation of the payload dynamics and degrade the overall transport performance. It is important to note that the simulation parameters were the same as those used in the previous subsection. The time evolution of the control variables is shown in Figure \ref{fig:Compensador_Variabes}, where the inclusion of the dynamic compensator results in a more precise and smoother speed response. In particular, the UAV is able to follow reference commands with reduced tracking error, avoiding the abrupt variations that can occur when only external loop compensation is used. As a result, closed-loop behavior becomes more regular, which also has a favorable indirect effect on payload movement, as the suspended mass is subjected to less aggressive action during the maneuver. From a quantitative perspective, this can be observed from the evolution of the speed tracking error and the oscillation-related performance indices (Figure \ref{fig:Compensador_Metricas}). The results show that incorporating the dynamic compensator reduces the magnitude of the internal loop tracking error throughout the maneuver, especially during the most demanding transient intervals. At the same time, the oscillation metrics remain bounded and exhibit a more favorable evolution than in the case without compensation, indicating that the UAV's improved response also translates into more stable payload movement.

In summary, the results obtained demonstrate that the dynamic compensator plays a fundamental role in the proposed control architecture. This can also be validated by analyzing the control errors (Figure \ref{fig:Compensador_ErroresControl}) and the speed tracking errors (Figure \ref{fig:Compensador_ErroresVelocidad}) between the actual speed and the references obtained from the kinematic controller. While the kinematic strategy shapes the reference motion to reduce oscillation, the dynamic compensator ensures that these references are effectively implemented at the UAV level. The combination of both layers results in improved closed-loop system performance, characterized by optimized speed tracking, smoother transients, and reduced payload oscillation during the transport task.


\begin{figure}
\centering
\includegraphics[width=0.95\linewidth]{Figures/Compensador.eps}
\caption{Time evolution of the reference and actual payload states under the proposed kinematic swing-compensation scheme. From top to bottom, the plots show the position-tracking responses in the $x$, $y$, and $z$ directions, along with the yaw-angle evolution. The results indicate accurate trajectory tracking with a substantial reduction in the oscillatory effects induced by the suspended payload.}\label{fig:Compensador_Variabes}
\end{figure}


\begin{figure}
\centering
\includegraphics[width=0.95\linewidth]{Figures/Compensador_Angles.eps}
\caption{Time evolution of the oscillation-related performance metrics obtained with the proposed kinematic swing-compensation scheme: payload swing angle $\alpha$, payload position error magnitude $\rho$, and tangential velocity norm $\|v_\perp\|$. The results show that, after the initial transient, the oscillatory motion remains bounded and significantly attenuated throughout the maneuver.}\label{fig:Compensador_Metricas}
\end{figure}

\begin{figure}
\centering
\includegraphics[width=0.95\linewidth]{Figures/Compensador_ErroresControl.eps}
\caption{Time evolution of the control errors.}\label{fig:Compensador_ErroresControl}
\end{figure}

\begin{figure}
\centering
\includegraphics[width=0.95\linewidth]{Figures/Compensador_ErroresVelocidad.eps}
\caption{Time evolution of the velocity errors.}\label{fig:Compensador_ErroresVelocidad}
\end{figure}




\subsection{Quantitative Performance Comparison}

Table \ref{tab:metrics} compares the metrics of simulations performed using only the kinematic controller and those performed using the proposed cascaded scheme (kinematic layer plus dynamic compensator). The metrics cover tracking quality (RMSE of payload position and velocity) and oscillation dampening ($\alpha$, $\rho$, and tangential velocity $v_T$). It can be observed that the addition of the dynamic compensator improves all reported indicators. Tracking errors decrease moderately (RMSE $\tilde{\mathbf{x}}$: $0.4350\!\to\!0.3671$; RMSE $\dot{\tilde{\mathbf{x}}}$: $0.3301\!\to\!0.3066$), while oscillation-related metrics show the greatest improvements, especially the maximum tangential velocity ($\max v_T$: $1.1420\!\to\!0.1403$, $\approx 88\%$). This result confirms the main advantage of the proposed architecture: the kinematic layer creates a payload-accounting reference, and the internal adaptive compensator ensures that the UAV actually tracks it, damping transient pendulum excitation without sacrificing trajectory tracking accuracy.
The observed 88\% reduction in tangential oscillation velocity demonstrates that the A-STA controller effectively restores performance symmetry, making the payload trajectory invariant to external wind gusts and cable state variations. This high level of stability is paramount for the safe transportation of sensitive cargo, ensuring that the system's operational envelope remains predictable even under asymmetric environmental disturbances.

\begin{table*}[t]
\centering
\caption{Performance metrics for the evaluated controller configurations. The $\boldsymbol{\Delta}$ row reports the relative improvement of the dynamic compensator with respect to the kinematic controller, computed as $(J_{\mathrm{Kin}}-J_{\mathrm{Dyn}})/J_{\mathrm{Kin}}\times100$ (positive values favor the dynamic compensator).}
\label{tab:metrics}
\renewcommand{\arraystretch}{1.15}
\setlength{\tabcolsep}{5pt}
\begin{tabular}{lcccccccc}
\toprule
Controller 
& RMSE $\dot{\tilde{\mathbf{x}}}$ 
& RMSE $\tilde{\mathbf{x}}$ 
& RMS$_{\alpha}$ 
& Max $\alpha$ 
& RMS$_{\rho}$ 
& Max $\rho$ 
& RMS$_{v_T}$ 
& Max $v_T$ \\
\midrule
Kinematic 
& 0.3301 & 0.4350 & 0.1582 & 0.3416 
& 0.3464 & 0.6902 & 0.2265 & 1.1420 \\
Dyn. Comp. 
& \textbf{0.3066} & \textbf{0.3671} & \textbf{0.1402} & \textbf{0.2869}
& \textbf{0.3203} & \textbf{0.5685} & \textbf{0.1265} & \textbf{0.1403} \\
\midrule
$\boldsymbol{\Delta}$ [\%]
& $+7.1$ & $+15.6$ & $+11.4$ & $+16.0$
& $+7.5$ & $+17.6$ & $+44.2$ & $+87.7$ \\
\bottomrule
\end{tabular}
\end{table*}

\subsection{Robustness Evaluation Under Wind Disturbances}
\label{sec:robustness_wind}

Unlike previous nominal scenarios, where the analysis focused primarily on tracking performance and oscillation reduction under standard operating conditions, this experiment was specifically designed to evaluate closed-loop behavior in a more demanding scenario. A three-dimensional lemniscate trajectory with vertical motion was selected. This reference trajectory provides a challenging scenario for evaluating the proposed controller, as it requires the payload to follow a cyclical lateral motion while simultaneously undergoing vertical displacement. The reference trajectory is defined as:
\begin{equation}
    \mathbf{x}_d =
    \begin{bmatrix}
    a\sin(ct)\cos(ct) \\
    b\sin(ct) \\
    6+d\sin(ct) \\
    \arctan\left(\dot{y}_d/\dot{x}_d\right)
    \end{bmatrix} ,
\end{equation}
where the parameters were selected as
\begin{equation}
a=10,\qquad b=7.5,\qquad c=0.025\,2\pi,\qquad d=1.
\end{equation}
In addition, a payload mass of 300 g was considered, and two wind disturbance intervals were introduced during the simulation. The purpose of including wind is to analyze the effect of external disturbances on the control loop and evaluate the ability of the proposed strategy to mitigate unwanted oscillatory motion induced by the suspended payload. These disturbances were introduced at different stages of the maneuver to evaluate the transient response of the closed-loop system and its ability to recover the desired behavior after each disturbance. The figure shows representative snapshots of the UAV-payload system during the nominal phase and under two wind-disturbance conditions. In particular, Fig.~\ref{fig:wind_three_views}(a) illustrates the system's response during the nominal range, while Fig.~\ref{fig:wind_three_views}(b) and Fig.~\ref{fig:wind_three_views}(c) show the system's behavior under lateral and oblique wind disturbances, respectively.
\begin{figure}[t]
\centering

\includegraphics[width=0.6\linewidth]{Figures/3D_Vista_10Sec.eps}\\
\small\textbf{(a)} Nominal phase

\vspace{4pt}

\includegraphics[width=0.6\linewidth]{Figures/3D_Vista_VientoLateral.eps}\\
\small\textbf{(b)} Lateral wind disturbance

\vspace{4pt}

\includegraphics[width=0.6\linewidth]{Figures/3D_Vista_Viento45.eps}\\
\small\textbf{(c)} Oblique wind disturbance

\caption{Simulation snapshots of the UAV--payload system under the lemniscate reference trajectory: \textbf{(a)} nominal phase; \textbf{(b)} lateral wind disturbance; and \textbf{(c)} oblique wind disturbance.}
\label{fig:wind_three_views}
\end{figure}
The temporal evolution of the payload's position and yaw angle is shown in Figure ~\ref{fig:Viento_Angles}. The results show that the payload can follow the desired reference trajectory with limited tracking errors, even in the face of wind disturbances. While transient deviations are observed during disturbance intervals, the controller is able to recover nominal tracking behavior without generating unstable oscillations or excessive payload movement. This behavior indicates that the proposed strategy maintains adequate trajectory tracking performance under external disturbances.
The oscillation-related performance variables are shown in Fig.~\ref{fig:Viento_ErroresControl}. These results make it possible to analyze the payload swing angle $\alpha$, the horizontal displacement-related error $\rho$, and the tangential velocity norm $\|v_\perp\|$. After the initial transient, the oscillatory motion remains bounded throughout the maneuver. Moreover, during the wind-disturbance intervals, the controller attenuates the induced swing and prevents the payload oscillations from growing persistently. This confirms that the proposed method preserves its swing-reduction capability even when the payload dynamics are excited by both the reference trajectory and external perturbations.
\begin{figure}
\centering
\includegraphics[width=1.0\linewidth]{Figures/Viento_ErroresControl.eps}
\caption{Time evolution of the reference and actual payload states under wind disturbances. From top to bottom, the plots show the position-tracking responses in the $x$, $y$, and $z$ directions, along with the yaw-angle evolution.}
\label{fig:Viento_ErroresControl}
\end{figure}
\begin{figure}
\centering
\includegraphics[width=1.0\linewidth]{Figures/Viento_Angles.eps}
\caption{
Time evolution of the oscillation-related performance metrics under wind disturbances: payload swing angle $\alpha$, payload displacement-related error $\rho$, and tangential velocity norm $\|v_\perp\|$.}
\label{fig:Viento_Angles}
\end{figure}
Finally, Fig.~\ref{fig:Viento_ErroresVelocidad} presents the time evolution of the velocity-tracking errors. These results show that the inner-loop control action remains effective despite aerodynamic disturbances. As shown in Figs.~\ref{fig:Viento_Angles} and~\ref{fig:Viento_ErroresControl}, although the wind gusts are applied from specific lateral and oblique directions (introducing a clear asymmetry), the controller manages to keep the tracking errors within similar bounds for both the x and y axes. This result proves that the proposed architecture successfully restores performance symmetry, as the system's ability to recover is not biased by the direction of the perturbation. 

\begin{figure}
\centering
\includegraphics[width=1.0\linewidth]{Figures/Viento_ErroresVelocidad.eps}
\caption{Time evolution of the velocity-tracking errors under wind disturbances.}
\label{fig:Viento_ErroresVelocidad}
\end{figure}
In summary, the results of this simulation demonstrate two key properties of the proposed approach. First, the controller exhibits robust closed-loop behavior, maintaining stable operation and acceptable path-tracking accuracy despite the presence of external wind disturbances. Second, the controller retains its oscillation reduction capability even when the payload motion is affected by a more aggressive reference path and repeated external disturbances. Therefore, this simulation provides a rigorous validation scenario in which both robustness and oscillation suppression performance are simultaneously evaluated.


\subsection{Benchmark Against a Published Baseline Controller}
\label{sec:benchmark_smc}

So far, the complete scheme has been proven successful. It was verified that the kinematic controller reduces load oscillations and that the dynamic compensator improves control error. However, to further validate the results, a comparison was made with an external reference. Therefore, the proposed controller is compared with a first-order sliding mode controller (SMC) baseline proposed in \cite{kusznir2020sliding} for cable-suspended load transport with a quadrotor.

First-order sliding mode control is a natural and widely documented reference for suspended load transport, as it shares the same family of sliding modes as the proposed time-based sliding mode control (STA) and directly answers the question of what is gained by moving from discontinuous first-order action to continuous second-order action. Unlike fully coupled geometric SE(3) designs or model-based passivity approaches, a first-order SMC does not require the full coupled UAV-cable-load dynamic model, can be implemented in the same standardized command interface used in this article (Eq.~\eqref{eq:ModeloLucas}), and is representative of a significant and well-cited family of controllers for overhead load transport. Therefore, the reference is reproducible, computationally comparable, and technically fair with respect to the proposed scheme.

To ensure a fair comparison that isolates the effect of the internal controller, both schemes share (i) the same high-fidelity simulator described at the beginning of Section ~\ref{sec:simulations}; (ii) the same upward circular payload reference trajectory; (iii) the same physical parameters listed in Table ~\ref{tab:sim_params}; (iv) the same initial condition; and (v) the same payload-oriented kinematic external loop with quasi-static oscillation compensation (Section ~\ref{sec:kinematic}). The only element that differs between the two implementations is the internal velocity tracking law. The proposed controller uses the composite adaptive supertorsion law defined in \eqref{eq:AccionControl}--\eqref{eq:ajuste2}, while the baseline uses a first-order sliding-mode action of the form
\begin{equation}
\mathbf{u}_{\mathrm{SMC}}
=
-\mathbf{K}_{\mathrm{SMC}}\,\dfrac{\mathbf{s}}{\|\mathbf{s}\|+\varphi},
\label{eq:smc_baseline}
\end{equation}
built on the same sliding surface $\mathbf{s}$ introduced in Section~\ref{sec:dynamic}, where $\mathbf{K}_{\mathrm{SMC}}\!>\!\mathbf{0}$ is a diagonal gain matrix, and $\varphi>0$ is a boundary-layer thickness used to attenuate chattering, following the classical practice of \cite{kusznir2020sliding}. This substitution isolates the contribution of the STA and its adaptation law from the effect of the outer kinematic compensation layer.

Table~\ref{tab:benchmark_smc} summarizes the tracking and oscillation-related metrics, together with two chattering diagnostics, computed over the same maneuver used in Section~\ref{sec:simulations}. Because both controllers share the payload-aware kinematic layer, the differences reported here are attributable exclusively to the inner control law. Three observations stand out. First, the proposed controller improves the payload position tracking, with a $6.7\%$ reduction in $\mathrm{RMSE}(\tilde{\mathbf{x}})$ ($0.4390\!\to\!0.4095$ m), while the payload-velocity tracking remains essentially equivalent for both schemes (a difference of $0.2\%$), which shows that the STA does not sacrifice velocity fidelity to attain the position gain. Second, all swing-related metrics ($\alpha$, $\rho$ and $v_T$) also favor the Super-Twisting controller, with a $3.4\%$ reduction in $\max v_T$, the metric most directly linked to the residual pendular energy of the load. Third, and consistent with the well-known chattering behavior of first-order sliding modes, the A-STA delivers a smoother command signal: the RMS derivative of the horizontal command $\mathrm{RMS}(\dot{u}_x)$ decreases by $11.1\%$ ($0.226\!\to\!0.201$ 1/s), and the total variation of $u_x$ decreases by $15.4\%$ ($14.615\!\to\!12.365$). On the $y$-axis, the two schemes produce nearly indistinguishable chattering diagnostics ($3\%\!-\!4\%$ apart), which is consistent with the fact that both share the same boundary-layer/second-order philosophy on the sliding surface; a systematic average across both horizontal axes still favors the A-STA in both diagnostics.
\begin{table}[t]
\centering
\caption{Benchmark against the first-order SMC baseline \cite{kusznir2020sliding} over the ascending-circular payload reference used in Section~\ref{sec:simulations}. Both controllers share the same simulator, trajectory, physical parameters, initial condition, and payload-aware kinematic outer loop; the only difference is the inner-loop action. The last four rows are chattering diagnostics computed on the normalized command signals sent to the vehicle. The $\boldsymbol{\Delta}$ column reports the relative improvement of the A-STA with respect to the SMC baseline (positive values favor the A-STA).}
\label{tab:benchmark_smc}
\renewcommand{\arraystretch}{1.15}
\setlength{\tabcolsep}{5pt}
\begin{tabular}{lccc}
\toprule
Metric & \shortstack{Baseline SMC\\Kusznir 2020~\cite{kusznir2020sliding}} & \shortstack{Proposed \\ A-STA} & $\boldsymbol{\Delta}$~[\%] \\
\midrule
RMSE $\tilde{\mathbf{x}}$ [m]         & 0.4390 & \textbf{0.4095} & $+6.7$ \\
RMSE $\dot{\tilde{\mathbf{x}}}$ [m/s] & 0.3132 & 0.3138          & $-0.2$ \\
RMS $\alpha$ [rad]                    & 0.1614 & \textbf{0.1593} & $+1.3$ \\
Max $\alpha$ [rad]                    & 0.3137 & \textbf{0.3124} & $+0.4$ \\
RMS $\rho$ [m]                        & 0.3226 & \textbf{0.3184} & $+1.3$ \\
Max $\rho$ [m]                        & 0.6203 & \textbf{0.6169} & $+0.6$ \\
RMS $v_T$ [m/s]                       & 0.1278 & \textbf{0.1265} & $+1.0$ \\
Max $v_T$ [m/s]                       & 0.1468 & \textbf{0.1418} & $+3.4$ \\
\midrule
$\mathrm{RMS}(\dot{u}_x)$ [1/s]       & 0.226  & \textbf{0.201}  & $+11.1$ \\
$\mathrm{RMS}(\dot{u}_y)$ [1/s]       & 0.217  & 0.225           & $-3.7$ \\
$\mathrm{Total Variation}(u_x)$ [-]                & 14.615 & \textbf{12.365} & $+15.4$ \\
$\mathrm{Total Variation}(u_y)$ [-]                & 18.169 & 18.722          & $-3.0$ \\
\bottomrule
\end{tabular}
\end{table}
Figures~\ref{fig:benchmark_metrics} and~\ref{fig:benchmark_control} illustrate the time-domain behavior underlying the numeric summary of Table~\ref{tab:benchmark_smc}. Figure~\ref{fig:benchmark_metrics} shows the swing angle $\alpha$, the horizontal displacement $\rho$, and the tangential velocity norm $\|v_\perp\|$ for both controllers over the same maneuver: both schemes are stable and payload-following, with envelopes of similar shape but a systematically smaller peak amplitude for the A-STA, particularly in $\|v_\perp\|$. Figure~\ref{fig:benchmark_control} shows the inner-loop correction signal along the horizontal axes. Along the $x$ axis, the first-order SMC exhibits a persistent low-amplitude sawtooth pattern that is smoothed but not eliminated by the boundary layer $\varphi$, whereas the A-STA action remains continuous with a smoother sinusoidal envelope. Along the $y$ axis, the STA additionally shows a large slow transient during the first ten seconds of the maneuver, corresponding to the online adjustment of the adaptive gains \eqref{eq:ajuste2} toward values consistent with the actual $\boldsymbol{\kappa}_u,\boldsymbol{\kappa}_v$; this excursion is confined to the adaptation transient and vanishes once the gains settle, which explains the slight $y$-axis chattering diagnostic in Table~\ref{tab:benchmark_smc} without contradicting the smoother steady-state behavior of the STA. Overall, this figure is fully consistent with the derivative-based chattering diagnostics reported in Table~\ref{tab:benchmark_smc}.
\begin{figure}[t]
\centering
\includegraphics[width=0.95\linewidth]{Figures/Comp_Baseline_Oscilaciones.eps}
\caption{Oscillation-related metrics for the proposed A-STA controller and the first-order SMC baseline \cite{kusznir2020sliding}: payload swing angle $\alpha$, horizontal displacement $\rho$, and tangential velocity norm $\|v_\perp\|$. Both controllers share the same payload-aware kinematic outer loop; the only difference is the inner-loop action.}
\label{fig:benchmark_metrics}
\end{figure}

\begin{figure}[t]
\centering
\includegraphics[width=0.95\linewidth]{Figures/Comp_Baseline_Chattering.eps}
\caption{Inner-loop correction signal along the horizontal axes for the proposed A-STA controller and the first-order SMC baseline. The higher-frequency activity of the SMC action illustrates the well-known chattering behavior of first-order sliding modes, which the super-twisting structure attenuates.}
\label{fig:benchmark_control}
\end{figure}
Two aspects of the comparison deserve emphasis. First, the absolute magnitude of the differences observed at the inner-loop level is intentionally modest: because both controllers share the payload-aware kinematic outer loop introduced in Section~\ref{sec:kinematic}, the largest oscillation reduction with respect to the purely reactive baseline was already reported in Section~\ref{sec:simulations} (see, in particular, the $\approx 88\%$ reduction in $\max v_T$ of Table~\ref{tab:metrics}). The role of the present benchmark is not to compete with that upper-level effect, but to isolate the gains from using a composite adaptive super-twisting action instead of a classical first-order sliding-mode action within the same outer loop. Second, the improvement is not concentrated in a single metric but appears consistently across the reported indices: the proposed A-STA controller matches the payload-velocity tracking of the SMC baseline, improves the payload-position RMSE by $6.7\%$, systematically reduces every swing-related metric (with the largest gain, $3.4\%$, on the operationally relevant peak tangential velocity), and reduces the horizontal-axis chattering diagnostics by $11$--$15\%$, without the need to select or retune a boundary-layer thickness $\varphi$. Beyond these numeric gains, the A-STA also carries a formal Lyapunov guarantee under bounded disturbance derivatives (Theorem~\ref{th:inner}) and an online adaptation law that would extend the same conclusions to payload masses or aerodynamic conditions different from those used in this benchmark, whereas the fixed-gain SMC would in general require re-tuning. This benchmark therefore complements the internal ablation of Section~\ref{sec:simulations} with an external, published reference and supports the practical relevance of the composite adaptive super-twisting architecture proposed in this paper.

\section{Conclusion}\label{sec:conclusiones}

This study developed a composite, payload-centric control architecture for transporting a cable-suspended load with a single UAV. The proposed design combines an outer kinematic loop that generates velocity references that account for oscillation with an adaptive, super-torsion inner loop that robustly tracks these references. Its main contribution lies in the coordinated operation of these two layers. Instead of treating every cable tilt as an unwanted oscillation, the outer loop considers the quasi-static tilt required by the commanded acceleration and compensates only for motion relative to that equilibrium. The inner loop then addresses parametric mismatch, unmodeled dynamics, and external disturbances. This modular formulation reduces the reliance on an exact model of the coupled UAV-cable-payload dynamics. It avoids constrained online optimization, making it a relatively lightweight alternative to payload-oriented trajectory tracking.

Stability analysis was performed separately for each layer. Under the idealized condition of perfect velocity tracking ($\dot{\mathbf{x}}=\mathbf{u}_k$), the Theorem~\ref{th:outer} establishes the input-to-state stability of the outer loop error cascade, with global exponential stability recovered in the unperturbed case. Independently, the Theorem~\ref{th:inner} shows, by composite Lyapunov analysis, that the velocity-tracking sliding variable and the parameter estimation error remain uniformly bounded under the coincidence condition set in $\mathbf{K}_2$, and that the velocity-tracking error converges asymptotically to zero, with parameter convergence requiring additional persistent excitation of the regressor. These two results are complementary, but not yet unified: the formal demonstration that interconnecting the two loops (where the inner loop achieves approximate, not exact, velocity tracking) preserves the stability margins established for the idealized outer loop (e.g., using a singular perturbation or ISS cascade argument on the actual closed-loop system) is left for future work. Therefore, the simulation results presented here should be interpreted as strong empirical support for this composition, complementing, rather than replacing, a fully formalized interconnection demonstration.

High-fidelity simulations support this description. The roll-sensitive kinematic layer stabilized a tracking task that the uncompensated formulation could not complete, and the addition of the dynamic compensator improved all reported metrics in the inner ablation, including an approximate 88\% reduction in the maximum tangential roll velocity. During the more demanding lemniscate maneuver, the controller maintained bounded tracking and roll responses under both lateral and oblique wind disturbances and recovered its nominal behavior after each disturbance interval. Comparison to the published first-order SMC baseline provides an additional external benchmark. With the simulator, trajectory, physical parameters, initial conditions, and payload-aware external loop fixed, the proposed A-STA reduced the RMSE of payload position by 6.7\%, while the RMSE of payload velocity remained virtually unchanged (-0.2\%). It also improved all reported oscillation metrics, with reductions ranging from 0.4\% to 3.4\%, the largest being for maximum tangential velocity. The command smoothness results were more direction-dependent: along the x-axis, the RMS command derivative and total variation decreased by 11.1\% and 15.4\%, respectively, while the corresponding y-axis measurements were approximately 3–4\% higher due to the initial adaptation transient. Therefore, external benchmarking indicates a modest but systematic improvement in payload tracking and oscillation damping compared to first-order sliding mode control, along with a qualified reduction in vibrations, rather than a uniform advantage across all control channels.

The immediate next step is therefore experimental validation. This effort will include implementation on an onboard PX4/ROS\,2 platform, reconstruction of $\mathbf{h}$ and $\dot{\mathbf{h}}$ from lightweight sensing, and explicit characterization of measurement noise, delay, computation time, and actuator saturation. Indoor motion-capture experiments with different payload masses will assess adaptation under genuine parametric variation, followed by outdoor tests under natural wind. The same tracking and swing metrics used here will enable direct comparison between simulation and flight data. Future studies should also examine more aggressive trajectories, elastic or variable-length cables, different payload geometries, and a broader benchmark set that includes geometric, passivity-based, and predictive controllers under matched experimental conditions.
%%%%%%%%%%%%%%%%%%%%%%%%%%%%%%%%%%%%%%%%%%%%%%%%%%%%%%%%%%%%%%%%%%%%%%%%%%%%%%%%%%%%%%%%%%%%%%%%%%%%%%%%%%%%%%%%%%%%%%%%%%%%%%%%%%%%%%%%%%%%%%%%%%%%%%%%%%%%%%%%%%%%%%%%%%%%%%%%%%%%%%%%%%%%%%%%%%%%%%%%%%%%%%%%%%%%%%%%%%%%%%%%%%%%%%%%%%%%%%%%%%%%%%%%%%%%%%%%%%%%%%%%%%%%%%%%%%%%%%%


\backmatter

\bmhead{Supplementary information}

If your article has accompanying supplementary file/s please state so here. 

Authors reporting data from electrophoretic gels and blots should supply the full unprocessed scans for key as part of their Supplementary information. This may be requested by the editorial team/s if it is missing.

Please refer to Journal-level guidance for any specific requirements.

\bmhead{Acknowledgements}

This work was funded by the University of Jeddah, Jeddah, Saudi Arabia, under grant No. (UJ-22-KAP2-20588-1.). Therefore, the authors thank the University of Jeddah for its technical and financial support.

\bmhead{Funding}

 University of Jeddah, Jeddah, Saudi Arabia - grant No. (UJ-22-KAP2-20588-1.). 

\section*{Declarations}

Some journals require declarations to be submitted in a standardised format. Please check the Instructions for Authors of the journal to which you are submitting to see if you need to complete this section. If yes, your manuscript must contain the following sections under the heading `Declarations':

\begin{itemize}
\item Data availability: All data generated or analyzed during this study are included in this published article. The contributions of the current study are theoretical and supported by numerical simulations. The codes for simulations are available from the corresponding link:
\href{https://github.com/cdrosales84/uav-payload-asta-validation}{Code},
\href{https://youtu.be/G8a1F97qs_E?si=gQIMWdCDfgZr-XQR}{Video}. 
\item University of Jeddah, Jeddah, Saudi Arabia, under grant No. (UJ-22-KAP2-20588-1.)
\item No conflicts of interest.
\item Ethics approval and consent to participate
\item Author contribution:  
Conceptualization, S.K., S.B., F.S.A. and F. R.; methodology, S.K., C.R., E. O. F. and F. R.; software, S. K.; validation, S.B.; formal analysis, F.R.; investigation, F. S. A.; resources, S.K.; writing—original draft preparation, C.R., S. B., E.O.F. and F. R.; writing—review and editing, F.S.A. and S. K.; visualization, C.R. and E.O.F. ; supervision, S.B.; project administration, S. B.; funding acquisition, F.S.A. All authors reviewed the results and approved the final version of the manuscript.
\end{itemize}

\noindent
If any of the sections are not relevant to your manuscript, please include the heading and write `Not applicable' for that section. 

%%===================================================%%
%% For presentation purpose, we have included        %%
%% \bigskip command. Please ignore this.             %%
%%===================================================%%
\bigskip
\begin{flushleft}%
Editorial Policies for:

\bigskip\noindent
Springer journals and proceedings: \url{https://www.springer.com/gp/editorial-policies}

\bigskip\noindent
Nature Portfolio journals: \url{https://www.nature.com/nature-research/editorial-policies}

\bigskip\noindent
\textit{Scientific Reports}: \url{https://www.nature.com/srep/journal-policies/editorial-policies}

\bigskip\noindent
BMC journals: \url{https://www.biomedcentral.com/getpublished/editorial-policies}
\end{flushleft}



%\bibliographystyle{sn-mathphys-num}
\bibliography{References}


%\newpage
%If you do not have or do not want to include a photo, you can use IEEEbiographynophoto as shown below:

\end{document}
